% ============================================================
%  FICHIER : chapitre3_conclusion.tex
%  Contenu : Suite Section 2.2 Conception + Chapitre 3 + Conclusion
%  Usage   : % ============================================================
%  FICHIER : chapitre3_conclusion.tex
%  Contenu : Suite Section 2.2 Conception + Chapitre 3 + Conclusion
%  Usage   : % ============================================================
%  FICHIER : chapitre3_conclusion.tex
%  Contenu : Suite Section 2.2 Conception + Chapitre 3 + Conclusion
%  Usage   : % ============================================================
%  FICHIER : chapitre3_conclusion.tex
%  Contenu : Suite Section 2.2 Conception + Chapitre 3 + Conclusion
%  Usage   : \input{chapitre3_conclusion} dans complet.tex,
%            juste après le contenu de la Section 2.2 existante.
% ============================================================

% ============================================================
%  SUITE — SECTION 2.2 CONCEPTION DU SYSTÈME
% ============================================================

\subsection{Architecture Technique}
\label{ssec:archi}

L'architecture retenue pour la plateforme EEC~Géolocalisation est une
\textbf{architecture à services découplés}, dans laquelle chaque composant
fonctionnel est isolé dans un conteneur indépendant et communique avec les
autres via des interfaces standardisées. Ce choix s'oppose à une architecture
monolithique dans laquelle toutes les fonctions seraient regroupées dans un
seul processus~: la séparation des responsabilités facilite la maintenance,
permet la mise à jour indépendante de chaque couche et renforce la sécurité en
cloisonnant les surfaces d'attaque.

La plateforme est décomposée en cinq services orchestrés par Docker~Compose,
représentés dans la Figure~\ref{fig:diag-archi}~:

\begin{itemize}[itemsep=3pt, topsep=5pt, leftmargin=1.5cm]
  \item \textbf{Backend} --- serveur Django~5 avec GeoDjango et
    Django REST Framework~: expose l'intégralité de la logique métier
    sous forme d'API REST, applique le contrôle d'accès basé sur les rôles
    et gère les sessions d'authentification par cookies sécurisés
    (\textit{HttpOnly}, \textit{SameSite=Strict})~;
  \item \textbf{Base de données} --- PostgreSQL~16 avec l'extension
    PostGIS~3.4~: assure la persistance des données institutionnelles et
    des géométries spatiales (points, polygones)~; les requêtes de
    proximité exploitent les fonctions de calcul de distance natives~;
  \item \textbf{Serveur cartographique} --- GeoServer~2.26~: publie les
    couches géographiques de l'EEC selon les standards OGC WMS~1.3.0 et
    WFS~2.0.0~; les styles visuels sont définis via des fichiers SLD aux
    couleurs institutionnelles~;
  \item \textbf{Cache} --- Redis~7~: stocke les sessions Django côté
    serveur et met en cache les réponses fréquentes de l'API pour réduire
    la charge sur la base de données~;
  \item \textbf{Frontend} --- Next.js~14 avec l'App Router~: sert
    l'interface publique (carte Leaflet, clustering, filtres) et
    l'interface d'administration multi-niveaux~; les pages critiques sont
    rendues côté serveur (\textit{SSR}) pour garantir la sécurité des
    données sensibles.
\end{itemize}

Les échanges entre le frontend et le backend transitent exclusivement par
l'API REST en JSON. Le frontend Leaflet interroge GeoServer directement
pour les couches cartographiques, sans passer par le backend Django~: cette
séparation permet d'exploiter les protocoles OGC standard sans surcharger
le serveur applicatif. L'ensemble des services partage un réseau Docker
interne isolé~; seuls les ports nécessaires sont exposés vers l'hôte.

\begin{figure}[H]
  \centering
  \fbox{\parbox{0.88\textwidth}{\centering\vspace{3.2cm}
    \textit{[Diagramme d'architecture technique à services découplés
    --- à intégrer après export draw.io]}%
    \vspace{3.2cm}}}
  \caption{Architecture technique à services découplés de la plateforme EEC}
  \label{fig:diag-archi}
\end{figure}

%------------------------------------------------------------
\subsection{Modèle de Données}
\label{ssec:modele-bd}

Le schéma de données de la plateforme est un modèle relationnel-spatial
conforme aux normes PostGIS. Il est organisé en trois groupes d'entités
directement dérivés du domaine institutionnel de l'EEC et des exigences
fonctionnelles identifiées à la section~\ref{ssec:exigences}.

\noindent\textbf{Entités géographiques hiérarchiques.}
La structure territoriale de l'EEC est modélisée par quatre tables liées
en cascade~: \texttt{RegionSynodale}, qui stocke le polygone géospatial
de chacune des 22~régions (type \texttt{MultiPolygon}, SRID~4326)~;
\texttt{District}, rattaché par clé étrangère à une région~; \texttt{Paroisse},
rattachée à un district, avec un champ \texttt{position} de type
\texttt{Point} (SRID~4326) pour les coordonnées GPS, un indicateur
d'activité et les informations de contact~; \texttt{OEuvre}, rattachée à
une paroisse, avec un champ de type d'établissement (école, centre médical,
exploitation agropastorale, immeuble) et un champ de position facultatif.

\noindent\textbf{Ressources humaines.}
L'entité \texttt{Ouvrier} modélise le personnel ecclésiastique~: pasteur,
évangéliste ou diacre permanent, avec leur fonction et leur affectation à
une paroisse. L'entité \texttt{StatistiqueAnnuelle} est associée à une
paroisse et à une année~; elle enregistre les indicateurs démographiques
et financiers annuels.

\noindent\textbf{Gestion de la plateforme.}
L'entité \texttt{CompteUtilisateur} étend le modèle d'utilisateur standard~;
elle porte le rôle, les clés étrangères de périmètre synodal (région,
district ou paroisse), un champ de permissions personnalisées et un
indicateur de première connexion. L'entité \texttt{LogActivite} enregistre
de façon immuable chaque opération d'écriture~: auteur, type d'action,
modèle cible, identifiant de l'objet et horodatage UTC. Les entités visiteur
(\texttt{ProfilVisiteur}, \texttt{ParoisseEnregistree},
\texttt{ItinerairePersonnel}) complètent le schéma pour le module public
authentifié.

\begin{figure}[H]
  \centering
  \fbox{\parbox{0.88\textwidth}{\centering\vspace{3.2cm}
    \textit{[Schéma relationnel-spatial PostGIS
    --- à intégrer après génération draw.io]}%
    \vspace{3.2cm}}}
  \caption{Schéma de la base de données relationnelle-spatiale PostGIS}
  \label{fig:diag-bd}
\end{figure}

%------------------------------------------------------------
\subsection{Diagrammes de Séquence Système}
\label{ssec:seq-systeme}

Les diagrammes de séquence système modélisent les interactions entre les
acteurs externes et le système considéré comme une boîte noire~: ils montrent
les messages échangés et leur ordre chronologique, sans révéler les composants
internes. Deux scénarios représentatifs sont retenus.

\subsubsection{Séquence Système~1~--- Authentification et Accès au Tableau de Bord}

Ce diagramme (Figure~\ref{fig:seq-sys1}) représente le flux nominal de
connexion d'un administrateur~: saisie des identifiants, vérification par le
système, ouverture de session et redirection vers le tableau de bord adapté
au rôle. Il inclut le scénario alternatif de première connexion qui déclenche
obligatoirement le changement de mot de passe avant l'accès aux fonctions
d'administration.

\begin{figure}[H]
  \centering
  \fbox{\parbox{0.88\textwidth}{\centering\vspace{3.5cm}
    \textit{[Diagramme de séquence système 1
    --- Authentification et accès au tableau de bord
    --- à intégrer après génération draw.io]}%
    \vspace{3.5cm}}}
  \caption{Diagramme de séquence système --- Authentification et accès au tableau de bord}
  \label{fig:seq-sys1}
\end{figure}

\subsubsection{Séquence Système~2~--- Consultation Cartographique avec Filtrage}

Ce diagramme (Figure~\ref{fig:seq-sys2}) représente le flux de consultation
de la carte interactive par un utilisateur~: chargement de la carte, affichage
des couches institutionnelles, application d'un filtre hiérarchique par région
et district, mise à jour dynamique des marqueurs et consultation de la fiche
d'une entité sélectionnée. Il illustre la coordination entre l'interface
cartographique et le serveur de données géographiques.

\begin{figure}[H]
  \centering
  \fbox{\parbox{0.88\textwidth}{\centering\vspace{3.5cm}
    \textit{[Diagramme de séquence système 2
    --- Consultation cartographique avec filtrage hiérarchique
    --- à intégrer après génération draw.io]}%
    \vspace{3.5cm}}}
  \caption{Diagramme de séquence système --- Consultation cartographique avec filtrage}
  \label{fig:seq-sys2}
\end{figure}

%------------------------------------------------------------
\subsection{Diagrammes de Séquence Technique}
\label{ssec:seq-technique}

Les diagrammes de séquence technique descendent au niveau des composants
internes du système~: ils montrent les appels entre le frontend, le backend,
la base de données et le serveur cartographique, avec les noms de méthodes
et les flux de données réels.

\subsubsection{Séquence Technique~1~--- Gestion de Session Django}

Ce diagramme (Figure~\ref{fig:seq-tech1}) détaille le traitement interne
d'une requête d'authentification~: la validation des identifiants par le
backend, la création d'une entrée de session dans Redis, la génération du
cookie sécurisé \textit{HttpOnly} et sa transmission au navigateur, ainsi
que la vérification du jeton CSRF sur les requêtes modifiantes suivantes.
Il illustre la chaîne de sécurité complète mise en œuvre entre le frontend
Next.js et le backend Django.

\begin{figure}[H]
  \centering
  \fbox{\parbox{0.88\textwidth}{\centering\vspace{3.5cm}
    \textit{[Diagramme de séquence technique 1
    --- Authentification Django avec gestion de session et CSRF
    --- à intégrer après génération draw.io]}%
    \vspace{3.5cm}}}
  \caption{Diagramme de séquence technique --- Authentification Django et gestion de session}
  \label{fig:seq-tech1}
\end{figure}

\subsubsection{Séquence Technique~2~--- Requête Cartographique Leaflet vers GeoServer}

Ce diagramme (Figure~\ref{fig:seq-tech2}) détaille la chaîne de traitement
d'une requête cartographique émise par Leaflet.js~: construction de la
requête WMS par le client, transmission à GeoServer, interrogation de
PostGIS pour les entités géographiques, application du style SLD, génération
de la tuile image PNG et rendu dans l'interface. Il illustre le rôle de
GeoServer comme couche intermédiaire entre la base de données spatiale et
le client cartographique.

\begin{figure}[H]
  \centering
  \fbox{\parbox{0.88\textwidth}{\centering\vspace{3.5cm}
    \textit{[Diagramme de séquence technique 2
    --- Requête WMS Leaflet → GeoServer → PostGIS
    --- à intégrer après génération draw.io]}%
    \vspace{3.5cm}}}
  \caption{Diagramme de séquence technique --- Requête cartographique WMS de Leaflet vers GeoServer}
  \label{fig:seq-tech2}
\end{figure}

%------------------------------------------------------------
\subsection{Diagramme de Classes Technique}
\label{ssec:classes-tech}

Le diagramme de classes technique (Figure~\ref{fig:diag-classes-tech})
représente le modèle d'implémentation de la plateforme~: il étend le
diagramme de classes métier (Figure~\ref{fig:diag-classes-metier}) avec
les détails d'implémentation propres à Django et à GeoDjango~:
types de données précis (\texttt{PointField}, \texttt{MultiPolygonField},
\texttt{JSONField}), visibilité des attributs, signatures des méthodes
(sérialiseurs DRF, méthodes de filtrage RBAC, helpers de calcul spatial).
Les dépendances entre les classes de modèles, les sérialiseurs et les vues
DRF sont représentées explicitement.

\begin{figure}[H]
  \centering
  \fbox{\parbox{0.88\textwidth}{\centering\vspace{3.2cm}
    \textit{[Diagramme de classes technique
    --- à intégrer après génération draw.io]}%
    \vspace{3.2cm}}}
  \caption{Diagramme de classes technique de l'implémentation Django/GeoDjango}
  \label{fig:diag-classes-tech}
\end{figure}

\medskip
\noindent\textbf{Bilan du chapitre.}
Ce chapitre a formalisé l'ensemble des besoins de la plateforme en
dix-huit exigences fonctionnelles et dix exigences non-fonctionnelles, puis
modélisé le système à travers six diagrammes d'analyse --- contexte,
packages, cas d'utilisation, classes métier et activité --- et cinq
diagrammes de conception --- architecture technique, schéma de données
PostGIS, deux diagrammes de séquence système, deux diagrammes de séquence
technique et diagramme de classes technique. Ces modèles constituent le
contrat de spécification sur lequel s'appuie l'implémentation détaillée au
\textbf{Chapitre~3}.

\clearpage

% ============================================================
%  CHAPITRE 3 — IMPLÉMENTATION ET RÉSULTATS
% ============================================================
\chapter{Implémentation et Résultats}
\label{chap:implementation}

\begin{chapplan}
  \item \textbf{Section~\ref{sec:outils} --- Outils et Technologies}~:
    justification des choix technologiques pour chaque composant de la
    pile logicielle~; tableau récapitulatif des outils, versions et rôles.
  \item \textbf{Section~\ref{sec:deploiement} --- Diagramme de Déploiement}~:
    distribution des cinq services sur les nœuds physiques et virtuels~;
    configuration réseau Docker Compose.
  \item \textbf{Section~\ref{sec:resultats} --- Résultats}~:
    présentation annotée des interfaces fonctionnelles~: carte interactive,
    authentification, tableaux de bord par rôle, import/export et journal
    d'audit~; confrontation aux exigences du Chapitre~2.
  \item \textbf{Section~\ref{sec:discussion} --- Discussion}~:
    vérification des objectifs, analyse critique des forces et des limites
    de la solution réalisée.
\end{chapplan}

\lettrine[lines=2, loversize=0.05]{L}{es} décisions de conception
formalisées au Chapitre~2 se concrétisent dans ce chapitre en une
implémentation fonctionnelle et déployable. La
section~\ref{sec:outils} justifie les choix technologiques opérés pour
chaque composant de la pile logicielle et dresse un tableau récapitulatif
des outils utilisés. La section~\ref{sec:deploiement} présente le diagramme
de déploiement UML qui décrit la distribution physique des services.
La section~\ref{sec:resultats} expose les résultats à travers les
interfaces fonctionnelles de la plateforme, commentées au regard des
exigences définies au Chapitre~2. La section~\ref{sec:discussion} propose
enfin une analyse critique des forces et des limites de la solution.

%------------------------------------------------------------
\section{Outils et Technologies}
\label{sec:outils}
%------------------------------------------------------------

\subsection{Justification des Choix Technologiques}
\label{ssec:justification}

Le choix de chaque technologie résulte d'une évaluation explicite des
alternatives disponibles au regard des contraintes identifiées~: contraintes
institutionnelles (déploiement sur infrastructure limitée, données
sensibles), contraintes fonctionnelles (requêtes spatiales, RBAC
hiérarchique, publication OGC) et contraintes de maintenabilité
(logiciels libres, documentation active, communauté de support).

\subsubsection{Django~5 et Django REST Framework}

Django~5~\cite{django2024} a été retenu comme framework backend principal
pour trois raisons déterminantes. En premier lieu, son extension spatiale
native GeoDjango offre une intégration directe avec PostGIS~: les requêtes
de distance (\texttt{Distance()}), les filtres spatiaux (\texttt{contains},
\texttt{intersects}) et les champs de géométrie (\texttt{PointField},
\texttt{MultiPolygonField}) sont disponibles via l'ORM sans SQL manuel.
En second lieu, le module d'authentification de Django~5 propose une gestion
de sessions robuste par cookies \textit{HttpOnly}~; ce mécanisme est
nativement protégé contre les attaques de type XSS, contrairement aux
jetons JWT transmis dans le stockage local du navigateur, qui présentaient
une vulnérabilité documentée dans la version antérieure de la plateforme.
En troisième lieu, l'écosystème Python disponible ---
\texttt{openpyxl} pour la lecture des fichiers Excel EEC,
\texttt{reportlab} pour la génération PDF, \texttt{django-cors-headers}
pour la politique d'origines croisées --- couvre l'ensemble des besoins
fonctionnels sans dépendances tierces lourdes. Django REST
Framework~\cite{drf2024} complète Django en exposant les ressources
institutionnelles sous forme d'API REST JSON consommée par le frontend.

\subsubsection{PostgreSQL~16 et PostGIS~3.4}

PostgreSQL~16 associé à PostGIS~3.4~\cite{postgis2023} constitue le
standard de facto pour les bases de données spatiales open source~: il est
utilisé dans des projets cartographiques institutionnels à l'échelle nationale
et internationale, dispose d'une indexation spatiale performante
(\textit{R-tree GiST}) et garantit la conformité aux spécifications OGC
Simple Features. La fonction \texttt{ST\_Distance()} native permet
d'identifier la paroisse la plus proche d'une position GPS en une seule
requête SQL, sans traitement applicatif. Aucun SGBD alternatif ne combine
ces fonctionnalités spatiales et une intégration aussi mature avec GeoDjango.

\subsubsection{GeoServer~2.26}

GeoServer~2.26~\cite{geoserver2024} a été retenu pour la publication des
couches cartographiques institutionnelles au regard de sa conformité native
aux protocoles OGC WMS~1.3.0 et WFS~2.0.0, qui garantissent
l'interopérabilité avec tout client cartographique standard. La
configuration des styles visuels via des fichiers SLD permet de définir les
couleurs institutionnelles de l'EEC (vert synodal, fond clair) directement
sur le serveur, sans traitement côté client. L'alternative d'une
publication de fichiers GeoJSON statiques, utilisée dans la version
antérieure, était incompatible avec la mise à jour dynamique des couches
et ne répondait pas aux exigences ENF05 d'interopérabilité OGC.

\subsubsection{Next.js~14}

Next.js~14~\cite{nextjs2024} avec l'App Router a été retenu pour le frontend
pour sa capacité à combiner rendu côté serveur (\textit{SSR}) et rendu côté
client (\textit{CSR}) dans la même application. Le \textit{SSR} est
utilisé pour les pages d'administration~: les données sensibles ne sont
jamais envoyées au navigateur avant vérification de la session, ce qui
réduit la surface d'attaque. Le rendu côté client est utilisé pour la carte
interactive, qui requiert des interactions dynamiques non compatibles avec
le \textit{SSR}. L'App Router de Next.js~14 permet de définir des
\textit{layouts} partagés par rôle d'administration (SUPER/REGION,
DISTRICT, PAROISSE), ce qui simplifie l'implémentation du RBAC côté
frontend.

\subsubsection{Leaflet.js et MarkerCluster}

Leaflet.js~\cite{leaflet2024} a été retenu pour la bibliothèque
cartographique cliente pour sa légèreté (environ 42~Ko minifié), sa licence
libre et l'étendue de son écosystème de plugins. Le plugin
\texttt{Leaflet.MarkerCluster} implémente le regroupement automatique des
565~marqueurs de paroisses en indicateurs numériques~: sans ce mécanisme,
la densité des marqueurs rendrait la carte illisible aux niveaux de zoom
nationaux. La prise en charge native des couches \textit{TileLayer.WMS}
permet de consommer directement les couches GeoServer sans adaptation.

\subsubsection{Redis~7 et Docker Compose}

Redis~7~\cite{redis2024} assure deux fonctions complémentaires~:
le stockage des sessions Django côté serveur (évitant leur persistance
en base de données relationnelle) et la mise en cache des réponses API
fréquentes. Docker Compose~\cite{docker2024} orchestre l'ensemble des
cinq services dans une configuration déclarative reproductible sur toute
machine hôte disposant de Docker~: c'est la réponse directe à l'exigence
ENF06 de portabilité et à la limite de déploiement identifiée dans la
plateforme GeoEEC~v1.0.

%-----
\subsection{Récapitulatif des Outils Utilisés}
\label{ssec:outils-table}

Le Tableau~\ref{tab:outils} dresse la liste complète des outils et
technologies utilisés dans le projet, avec leur version exacte et leur
rôle précis.

\begin{table}[H]
\centering
\caption{Récapitulatif des outils et technologies de la plateforme EEC}
\label{tab:outils}
\setlength{\tabcolsep}{5pt}
\renewcommand{\arraystretch}{1.22}
\resizebox{\linewidth}{!}{%
\begin{tabular}{%
  >{\bfseries}m{3.6cm}
  >{\centering\arraybackslash}m{1.6cm}
  m{9.0cm}}
\toprule
\rowcolor{eecgreen}
\multicolumn{1}{>{\bfseries\color{white}}m{3.6cm}}{Outil / Technologie} &
\multicolumn{1}{>{\bfseries\color{white}\centering\arraybackslash}m{1.6cm}}{Version} &
\multicolumn{1}{>{\bfseries\color{white}}m{9.0cm}}{Rôle dans le projet} \\
\midrule
\rowcolor{eeclightgreen!60}
Python           & 3.12   & Langage principal du backend \\
Django           & 5.0    & Framework web backend~; ORM~; middleware RBAC~; sessions \\
\rowcolor{eeclightgreen!60}
GeoDjango        & 5.0    & Extension spatiale de Django~; interface avec PostGIS \\
Django REST Framework & 3.15 & Construction et sérialisation des endpoints API REST \\
\rowcolor{eeclightgreen!60}
PostgreSQL       & 16     & Système de gestion de base de données relationnelle \\
PostGIS          & 3.4    & Extension spatiale~; types géométriques~; requêtes de distance \\
\rowcolor{eeclightgreen!60}
GeoServer        & 2.26   & Serveur cartographique OGC~; publication WMS et WFS~; styles SLD \\
Redis            & 7      & Cache serveur~; stockage des sessions Django \\
\rowcolor{eeclightgreen!60}
Next.js          & 14     & Framework frontend React~; App Router~; SSR et CSR \\
TypeScript       & 5.x    & Typage statique du code frontend \\
\rowcolor{eeclightgreen!60}
Tailwind CSS     & 3.x    & Framework de mise en forme CSS utilitaire \\
Leaflet.js       & 1.9    & Bibliothèque cartographique JavaScript interactive \\
\rowcolor{eeclightgreen!60}
Leaflet.MarkerCluster & 1.5 & Regroupement automatique des marqueurs sur la carte \\
openpyxl         & 3.1    & Lecture et écriture des fichiers Excel (.xlsx) \\
\rowcolor{eeclightgreen!60}
reportlab        & 4.x    & Génération des documents PDF (exports, rapports) \\
Docker           & 24     & Conteneurisation des services de la plateforme \\
\rowcolor{eeclightgreen!60}
Docker Compose   & 2.x    & Orchestration des cinq services en environnement multi-conteneurs \\
draw.io / diagrams.net & --- & Conception des diagrammes UML du rapport \\
\bottomrule
\end{tabular}}
\end{table}

%------------------------------------------------------------
\section{Diagramme de Déploiement}
\label{sec:deploiement}
%------------------------------------------------------------

Le diagramme de déploiement (Figure~\ref{fig:diag-deploiement}) représente
la distribution physique et virtuelle des composants logiciels de la
plateforme sur les nœuds d'exécution. Deux nœuds principaux sont identifiés~:
le \textbf{serveur institutionnel} et le \textbf{navigateur client}.

Sur le serveur institutionnel, Docker Compose crée un réseau interne
(\texttt{eec\_network}) dans lequel les cinq conteneurs communiquent~: le
conteneur \texttt{backend} (Django~5, port~8000) expose l'API REST au
frontend et interroge \texttt{db} (PostgreSQL/PostGIS, port~5432) et
\texttt{cache} (Redis, port~6379)~; le conteneur \texttt{geoserver}
(port~8080) accède à \texttt{db} pour lire les données géographiques et
répond aux requêtes WMS/WFS~; le conteneur \texttt{frontend} (Next.js~14,
port~3000) sert l'interface web. Vers l'extérieur, seul le port~3000
(interface) est exposé~; les services internes restent cloisonnés dans
le réseau Docker. Les volumes persistants (\texttt{postgres\_data},
\texttt{geoserver\_data}) garantissent la conservation des données entre
les redémarrages de conteneurs.

Sur le navigateur client, Leaflet.js charge les tuiles de fond de carte
depuis OpenStreetMap et interroge directement GeoServer (port~8080) pour
les couches institutionnelles WMS/WFS~; les requêtes API REST transitent
vers le backend via le frontend Next.js.

\begin{figure}[H]
  \centering
  \fbox{\parbox{0.88\textwidth}{\centering\vspace{3.2cm}
    \textit{[Diagramme de déploiement UML
    --- à intégrer après génération draw.io]}%
    \vspace{3.2cm}}}
  \caption{Diagramme de déploiement UML --- Distribution des services Docker Compose}
  \label{fig:diag-deploiement}
\end{figure}

%------------------------------------------------------------
\section{Résultats}
\label{sec:resultats}
%------------------------------------------------------------

Cette section présente les interfaces fonctionnelles de la plateforme
implémentée, commentées au regard des exigences définies au
Chapitre~2. Les captures d'écran sont organisées selon les quatre
grands modules~: interface publique cartographique, authentification,
administration multi-niveaux et gestion des données.

%-----
\subsection{Interface Publique Cartographique}
\label{ssec:res-carte}

L'interface publique est accessible sans authentification depuis la racine
de la plateforme. Elle répond directement aux exigences EF03, EF04, EF05
et EF16 relatives à la cartographie interactive.

La Figure~\ref{fig:res-carte-publique} illustre la carte au niveau de zoom
national~: les 553~paroisses importées sont regroupées en indicateurs de
cluster numériques selon leur densité géographique. Le fond de carte
OpenStreetMap fournit le contexte géographique camerounais~; les couches
institutionnelles (polygones des régions synodales, marqueurs des paroisses
et des œuvres) sont servies par GeoServer en WMS. La barre de filtres
hiérarchiques en haut de la carte permet de sélectionner une région synodale,
un district et un type d'entité~; la carte se met à jour dynamiquement
sans rechargement de page.

\begin{figure}[H]
  \centering
  % ===== IMAGE 1 : Carte publique avec clustering =====
  \fbox{\parbox{0.88\textwidth}{\centering\vspace{3.8cm}
    \textit{[Capture d'écran 1 --- Carte interactive publique avec clustering de marqueurs]}%
    \vspace{3.8cm}}}
  \caption{Interface publique --- Carte interactive avec regroupement automatique des marqueurs}
  \label{fig:res-carte-publique}
\end{figure}

La Figure~\ref{fig:res-popup} montre le panneau de détail qui s'affiche lors
d'un clic sur un marqueur individuel~: nom de la paroisse, district et région
d'appartenance, coordonnées GPS, contacts et statut d'activité. Cette fiche
contextuelle répond à l'exigence EF06 et est accessible à tout utilisateur
sans connexion.

\begin{figure}[H]
  \centering
  % ===== IMAGE 2 : Popup de détail d'une paroisse =====
  \fbox{\parbox{0.88\textwidth}{\centering\vspace{3.8cm}
    \textit{[Capture d'écran 2 --- Popup de détail d'une paroisse sur la carte]}%
    \vspace{3.8cm}}}
  \caption{Interface publique --- Popup de détail d'une paroisse sélectionnée}
  \label{fig:res-popup}
\end{figure}

%-----
\subsection{Interface d'Authentification}
\label{ssec:res-auth}

La Figure~\ref{fig:res-login} présente la page de connexion de l'espace
d'administration (\texttt{/admin/login}). Le formulaire soumet les
identifiants à l'API backend~; en cas de succès, un cookie de session
\textit{HttpOnly} est transmis et l'utilisateur est redirigé vers le
tableau de bord correspondant à son rôle. En cas d'identifiants incorrects,
un message d'erreur est affiché sans révéler si l'identifiant ou le mot de
passe est la source de l'erreur, conformément aux bonnes pratiques de
sécurité recommandées par l'OWASP~\cite{owasp2021}. Pour un compte dont
l'indicateur de première connexion est actif, la redirection s'effectue
automatiquement vers le formulaire de changement de mot de passe.

\begin{figure}[H]
  \centering
  % ===== IMAGE 3 : Page de connexion administrateur =====
  \fbox{\parbox{0.88\textwidth}{\centering\vspace{3.8cm}
    \textit{[Capture d'écran 3 --- Page de connexion de l'espace administrateur]}%
    \vspace{3.8cm}}}
  \caption{Interface d'authentification --- Page de connexion sécurisée}
  \label{fig:res-login}
\end{figure}

%-----
\subsection{Interface d'Administration Multi-Niveaux}
\label{ssec:res-admin}

L'interface d'administration adapte son contenu et ses fonctionnalités au
rôle de l'utilisateur connecté, conformément aux exigences EF02 et EF07 à
EF09. Quatre interfaces distinctes sont générées selon le rôle~: SUPER/REGION,
DISTRICT et PAROISSE.

La Figure~\ref{fig:res-dashboard} illustre le tableau de bord de
l'administrateur national (rôle SUPER)~: il affiche les indicateurs de
synthèse nationaux (nombre total de paroisses actives, de districts, de
régions, d'ouvriers et d'œuvres), les statistiques de couverture GPS par
région et l'évolution des indicateurs sur les dernières années. Les
indicateurs sont calculés dynamiquement par le backend en filtrant les
données selon le périmètre du rôle~: un administrateur régional voit les
mêmes indicateurs mais limités à sa région synodale.

\begin{figure}[H]
  \centering
  % ===== IMAGE 4 : Tableau de bord administrateur national =====
  \fbox{\parbox{0.88\textwidth}{\centering\vspace{3.8cm}
    \textit{[Capture d'écran 4 --- Tableau de bord de l'administrateur national]}%
    \vspace{3.8cm}}}
  \caption{Interface d'administration --- Tableau de bord de l'administrateur national}
  \label{fig:res-dashboard}
\end{figure}

La Figure~\ref{fig:res-paroisses} présente l'interface de gestion des
paroisses~: liste paginée filtrée par région et district, formulaire de
création et de modification avec validation des coordonnées GPS en temps
réel, et action de suppression avec confirmation. Le système de contrôle
d'accès RBAC restreint les opérations à la liste des paroisses appartenant
au périmètre synodal de l'utilisateur~: un administrateur de district ne
peut visualiser ni modifier les paroisses d'un autre district.

\begin{figure}[H]
  \centering
  % ===== IMAGE 5 : Interface de gestion des paroisses =====
  \fbox{\parbox{0.88\textwidth}{\centering\vspace{3.8cm}
    \textit{[Capture d'écran 5 --- Interface de gestion des paroisses (liste et formulaire)]}%
    \vspace{3.8cm}}}
  \caption{Interface d'administration --- Gestion des paroisses avec contrôle RBAC}
  \label{fig:res-paroisses}
\end{figure}

%-----
\subsection{Module d'Import et d'Export}
\label{ssec:res-import}

Le module d'import (Figure~\ref{fig:res-import}) répond à l'exigence EF10~:
il permet de charger les fichiers Excel fournis par les services de l'EEC
directement dans la base de données, après une validation ligne par ligne.
Le rapport d'import affiché après traitement détaille le nombre de lignes
insérées avec succès, le nombre de lignes rejetées et, pour chaque
rejet, la raison précise (coordonnées GPS hors territoire, district
inexistant, doublon détecté). Cette traçabilité de l'import répond
directement au problème de qualité des données identifié lors de l'audit
préliminaire~: 127~paroisses sources sans coordonnées GPS valides ont ainsi
été identifiées et signalées lors de la première importation.

\begin{figure}[H]
  \centering
  % ===== IMAGE 6 : Module d'import Excel et rapport de résultats =====
  \fbox{\parbox{0.88\textwidth}{\centering\vspace{3.8cm}
    \textit{[Capture d'écran 6 --- Module d'import Excel avec rapport détaillé de résultats]}%
    \vspace{3.8cm}}}
  \caption{Module d'import --- Chargement de données Excel et rapport de validation}
  \label{fig:res-import}
\end{figure}

%-----
\subsection{Journal d'Audit}
\label{ssec:res-audit}

Le journal d'audit (Figure~\ref{fig:res-audit}) est accessible
exclusivement à l'administrateur national. Il liste toutes les opérations
d'écriture effectuées sur les données institutionnelles, dans l'ordre
chronologique inversé~: création, modification ou suppression d'une entité,
avec l'auteur de l'action, le type d'entité concernée, l'identifiant de
l'objet et l'horodatage UTC. Des filtres par type d'action, type d'entité
et période permettent de cibler une recherche. En cliquant sur une entrée,
l'administrateur accède au détail de la modification~: valeurs avant et
après la mise à jour. Ce module répond aux exigences EF14 et EF15 relatives
à la traçabilité, et constitue l'un des apports distinctifs par rapport
à la plateforme GeoEEC~v1.0 qui ne disposait d'aucun mécanisme d'audit.

\begin{figure}[H]
  \centering
  % ===== IMAGE 7 : Journal d'audit =====
  \fbox{\parbox{0.88\textwidth}{\centering\vspace{3.8cm}
    \textit{[Capture d'écran 7 --- Journal d'audit des modifications institutionnelles]}%
    \vspace{3.8cm}}}
  \caption{Journal d'audit --- Historique des opérations d'écriture sur les données}
  \label{fig:res-audit}
\end{figure}

%------------------------------------------------------------
\section{Discussion}
\label{sec:discussion}
%------------------------------------------------------------

\subsection{Adéquation aux Objectifs}
\label{ssec:disc-objectifs}

Les six objectifs formulés en Introduction Générale sont confrontés aux
résultats obtenus. L'objectif~1 (modélisation du schéma de données
hiérarchique) est atteint~: le schéma PostGIS implémenté couvre les
22~régions synodales, 137~districts, 553~paroisses importées, 311~œuvres
et 685~ouvriers, avec toutes les contraintes d'intégrité référentielle et
les champs de géolocalisation. L'objectif~2 (API REST sécurisée avec RBAC)
est atteint~: l'API expose 47~endpoints documentés, protégés par le
middleware RBAC à quatre  niveaux et par le journal d'audit automatique.
L'objectif~3 (configuration GeoServer OGC) est atteint~: trois couches
institutionnelles sont publiées en WMS~1.3.0 et WFS~2.0.0 avec styles SLD.
L'objectif~4 (interface cartographique Leaflet) est atteint~: la carte
interactive intègre le clustering automatique, les filtres hiérarchiques
et les popups de détail. L'objectif~5 (module d'administration multi-niveaux)
est atteint~: les quatre interfaces de rôle, les modules d'import Excel et
d'export PDF/Excel sont fonctionnels. L'objectif~6 (conteneurisation Docker)
est atteint~: la plateforme est entièrement déployable avec une commande
unique sur tout serveur disposant de Docker.

%-----
\subsection{Forces de la Solution}
\label{ssec:disc-forces}

\textbf{Sécurité améliorée par rapport à la version antérieure.}
Le remplacement de l'authentification JWT par des sessions Django avec
cookies \textit{HttpOnly} et \textit{SameSite=Strict} élimine le vecteur
d'attaque XSS qui permettait la capture des jetons transmis dans le stockage
local du navigateur. La protection CSRF systématique sur toutes les
requêtes modifiantes renforce ce dispositif, conformément aux recommandations
de l'OWASP~\cite{owasp2021}.

\textbf{Interopérabilité cartographique.}
La publication des couches géographiques via GeoServer en WMS~1.3.0 et
WFS~2.0.0 rend les données de l'EEC consommables par tout client
cartographique compatible OGC~: QGIS, OpenLayers, ou un futur portail
institutionnel camerounais. Aucune solution concurrente analysée au
Chapitre~1 ne proposait cette combinaison.

\textbf{RBAC aligné sur la hiérarchie réelle de l'EEC.}
Le modèle de contrôle d'accès à quatre niveaux --- national, régional,
district et paroissial --- correspond exactement à la structure synodale
de l'institution~: un administrateur de district peut gérer les paroisses
de son périmètre sans accès aux données des districts voisins. Ce niveau
de granularité était absent de la version antérieure.

\textbf{Traçabilité complète et automatique.}
Tout enregistrement de modification, de création ou de suppression
déclenche automatiquement l'écriture d'une entrée dans le journal d'audit
via les signaux Django, sans intervention du développeur dans chaque vue.
Cette automatisation garantit l'exhaustivité de l'historique, impossible
à assurer manuellement.

\textbf{Déploiement reproductible.}
La configuration Docker Compose en cinq services garantit que l'environnement
d'exécution est identique en développement et en production~: plus aucune
divergence de versions de bibliothèques ou de configuration système ne peut
causer des défaillances à l'installation.

Ces cinq forces confirment la pertinence des décisions architecturales
formalisées au Chapitre~2. Elles positionnent la plateforme de façon
distincte par rapport aux solutions analysées dans l'état de l'art, en
particulier sur les axes de sécurité, d'interopérabilité OGC et de
RBAC institutionnel. La section suivante expose les limites structurelles
qui tempèrent la portée de ces résultats.

%------------------------------------------------------------
\section{Limites et Contraintes}
\label{sec:limites}
%------------------------------------------------------------

Tout travail d'ingénierie comporte des contraintes qui en délimitent la
portée. Une évaluation honnête de ces limites est indispensable pour
apprécier correctement la valeur de la solution et orienter les
développements futurs.

\textbf{Couverture GPS incomplète.}
L'audit des données sources a révélé que 127~paroisses sur 565, soit
22~\% de l'effectif total, ne disposent d'aucune coordonnée GPS valide
dans les fichiers Excel fournis par l'EEC. Ces paroisses sont importées
en base de données mais ne s'affichent pas sur la carte interactive.
La plateforme ne peut donc pas prétendre à une couverture cartographique
complète du réseau ecclésiastique~; son utilité pour les décisions de
planification territoriale nationale reste limitée tant que ce taux
n'est pas réduit par une campagne de collecte terrain.

\textbf{Distance à vol d'oiseau.}
La fonction de calcul d'itinéraire repose sur la formule de Haversine,
qui mesure la distance géodésique en ligne droite entre deux points.
Cette mesure sous-estime systématiquement la durée réelle de trajet
sur un réseau routier camerounais caractérisé par des voies sinueuses
et un relief accidenté~; elle ne tient pas compte des barrières
physiques (fleuves, massifs montagneux, absence de pont).
La durée estimée à 60~km/h de moyenne est donc approximative et ne
peut servir de base à des décisions opérationnelles précises.

\textbf{Absence d'application mobile.}
La plateforme est conçue pour des écrans desktop et tablette
(résolution minimale 1\,024~$\times$~768~px, ENF09). Elle n'est pas
adaptée aux smartphones, qui constituent pourtant le principal outil
numérique des administrateurs paroissiaux en zone rurale camerounaise.
En l'absence d'application mobile dédiée, la collecte terrain des
coordonnées GPS des 127~paroisses non géolocalisées reste sans
outillage numérique adapté.

\textbf{Authentification à facteur unique.}
L'accès à l'espace d'administration repose sur un couple
identifiant/mot de passe sans mécanisme de double authentification
(TOTP ou code temporaire par messagerie). Pour des comptes de niveau
national disposant d'un accès complet à l'ensemble des données de
l'EEC, cette limite expose la plateforme aux risques de compromission
par capture ou réutilisation de mot de passe, en l'absence de second
facteur de vérification.

\textbf{Déploiement mono-instance.}
L'architecture Docker Compose actuelle déploie chaque service en
instance unique sur un seul serveur physique. Elle ne prévoit pas de
mécanisme de réplication ou de bascule automatique en cas de défaillance
matérielle~; la disponibilité de la plateforme est donc directement
dépendante de la fiabilité du serveur institutionnel hôte.

\medskip
\noindent\textbf{Bilan du chapitre.}
Ce chapitre a présenté et justifié les choix technologiques de la
plateforme à travers quatre tableaux comparatifs, exposé les résultats
à travers sept interfaces fonctionnelles commentées, analysé les cinq
forces distinctives de la solution dans la section Discussion, puis
identifié cinq limites structurelles dans cette section. La
\textbf{Conclusion Générale} synthétise l'ensemble de la démarche
et formule les perspectives d'amélioration directement dérivées
de ces limites.

\clearpage

% ============================================================
%  CONCLUSION GÉNÉRALE
% ============================================================
\chapter*{Conclusion Générale}
\addcontentsline{toc}{chapter}{Conclusion Générale}
\markboth{Conclusion Générale}{Conclusion Générale}

\noindent
L'Église Évangélique du Cameroun gérait jusqu'alors ses 565~paroisses,
311~œuvres sociales et 685~ouvriers ecclésiastiques répartis sur
22~régions synodales et 137~districts sans aucun outil numérique
centralisé~: pas de visualisation géographique, pas de contrôle d'accès
formalisé selon les périmètres synodaux, pas de journal d'audit des
modifications, pas de publication cartographique interopérable. Ce travail
a abordé ce problème par la conception et le développement d'une
plateforme web cartographique institutionnelle, en s'appuyant sur
l'hypothèse qu'une architecture à services découplés, entièrement
conteneurisée, combinant un serveur cartographique OGC et un contrôle
d'accès hiérarchique aligné sur la structure synodale, permettrait d'y
répondre de façon complète et déployable dans un contexte institutionnel
à ressources limitées.

Cette hypothèse est validée par les résultats obtenus~: la plateforme
implémentée déploie en une seule commande cinq services
--- Django~5~+~GeoDjango, PostgreSQL~+~PostGIS~3.4, GeoServer~2.26,
Next.js~14 et Redis~7~--- et offre une carte interactive avec clustering
de 553~paroisses, des filtres hiérarchiques par région et district, un
tableau de bord adaptatif par rôle, un module d'import/export des données
institutionnelles, et un journal d'audit automatique et exhaustif de toutes
les opérations d'écriture. L'ensemble des six objectifs formulés en
Introduction Générale a été atteint. Par rapport à la plateforme
GeoEEC~v1.0, la solution apporte une valeur ajoutée précise et documentée
sur cinq axes~: sécurité de l'authentification par sessions \textit{HttpOnly},
RBAC aligné sur les quatre niveaux de la hiérarchie synodale réelle,
publication cartographique conforme aux standards OGC WMS et WFS, traçabilité
automatique par journal d'audit, et reproductibilité de l'environnement
d'exécution par conteneurisation complète.

\medskip
\noindent\textbf{Limites.}
Cinq limites structurelles, analysées en détail dans la
section~\ref{sec:limites} du Chapitre~3, tempèrent la portée de ces
résultats. La couverture cartographique reste incomplète~: 127~paroisses
(22~\%) ne disposent d'aucune coordonnée GPS dans les fichiers sources de
l'EEC et ne s'affichent pas sur la carte, ce qui restreint directement
l'utilité de la plateforme pour les décisions de planification territoriale.
Le calcul d'itinéraire repose sur la distance à vol d'oiseau et non sur
le réseau routier réel, produisant des durées de trajet sous-estimées dans
les zones à relief accidenté. L'absence d'interface mobile prive les
administrateurs paroissiaux d'un outil de saisie GPS sur le terrain.
L'authentification à facteur unique expose les comptes nationaux à des
risques de compromission par mot de passe seul. Enfin, le déploiement
mono-instance ne prévoit pas de mécanisme de haute disponibilité en cas
de défaillance du serveur hôte.

\medskip
\noindent\textbf{Perspectives.}
Ces limites définissent directement les axes de développement prioritaires.
La première perspective est le développement d'une application mobile
légère (React Native ou Progressive Web App) permettant aux administrateurs
paroissiaux de saisir ou de corriger les coordonnées GPS directement sur le
terrain~: cette seule fonctionnalité permettrait de résorber le déficit de
géolocalisation des 127~paroisses non couvertes et de faire passer le taux
de couverture cartographique de 78~\% à 100~\%. La deuxième perspective est
l'intégration d'un moteur de calcul d'itinéraires routiers~--- via l'API
OSRM (\textit{Open Source Routing Machine}) ou une instance locale de
GraphHopper~--- pour remplacer la distance à vol d'oiseau par des temps de
trajet réalistes sur les routes camerounaises. La troisième perspective est
l'activation de la double authentification par TOTP (\textit{Time-based
One-Time Password}) pour les comptes d'administrateurs nationaux et
régionaux, en s'appuyant sur la bibliothèque \texttt{django-otp}
déjà compatible avec l'infrastructure Django. Enfin, la réintégration
du module de questions-réponses en langage naturel développé dans
GeoEEC~v1.0 --- enrichi des nouvelles entités et des statistiques annuelles
--- ouvrirait la plateforme à des usages analytiques avancés permettant
aux responsables synodaux d'interroger les données institutionnelles sans
compétences techniques.

\noindent\rule{\linewidth}{0.8pt}
 dans complet.tex,
%            juste après le contenu de la Section 2.2 existante.
% ============================================================

% ============================================================
%  SUITE — SECTION 2.2 CONCEPTION DU SYSTÈME
% ============================================================

\subsection{Architecture Technique}
\label{ssec:archi}

L'architecture retenue pour la plateforme EEC~Géolocalisation est une
\textbf{architecture à services découplés}, dans laquelle chaque composant
fonctionnel est isolé dans un conteneur indépendant et communique avec les
autres via des interfaces standardisées. Ce choix s'oppose à une architecture
monolithique dans laquelle toutes les fonctions seraient regroupées dans un
seul processus~: la séparation des responsabilités facilite la maintenance,
permet la mise à jour indépendante de chaque couche et renforce la sécurité en
cloisonnant les surfaces d'attaque.

La plateforme est décomposée en cinq services orchestrés par Docker~Compose,
représentés dans la Figure~\ref{fig:diag-archi}~:

\begin{itemize}[itemsep=3pt, topsep=5pt, leftmargin=1.5cm]
  \item \textbf{Backend} --- serveur Django~5 avec GeoDjango et
    Django REST Framework~: expose l'intégralité de la logique métier
    sous forme d'API REST, applique le contrôle d'accès basé sur les rôles
    et gère les sessions d'authentification par cookies sécurisés
    (\textit{HttpOnly}, \textit{SameSite=Strict})~;
  \item \textbf{Base de données} --- PostgreSQL~16 avec l'extension
    PostGIS~3.4~: assure la persistance des données institutionnelles et
    des géométries spatiales (points, polygones)~; les requêtes de
    proximité exploitent les fonctions de calcul de distance natives~;
  \item \textbf{Serveur cartographique} --- GeoServer~2.26~: publie les
    couches géographiques de l'EEC selon les standards OGC WMS~1.3.0 et
    WFS~2.0.0~; les styles visuels sont définis via des fichiers SLD aux
    couleurs institutionnelles~;
  \item \textbf{Cache} --- Redis~7~: stocke les sessions Django côté
    serveur et met en cache les réponses fréquentes de l'API pour réduire
    la charge sur la base de données~;
  \item \textbf{Frontend} --- Next.js~14 avec l'App Router~: sert
    l'interface publique (carte Leaflet, clustering, filtres) et
    l'interface d'administration multi-niveaux~; les pages critiques sont
    rendues côté serveur (\textit{SSR}) pour garantir la sécurité des
    données sensibles.
\end{itemize}

Les échanges entre le frontend et le backend transitent exclusivement par
l'API REST en JSON. Le frontend Leaflet interroge GeoServer directement
pour les couches cartographiques, sans passer par le backend Django~: cette
séparation permet d'exploiter les protocoles OGC standard sans surcharger
le serveur applicatif. L'ensemble des services partage un réseau Docker
interne isolé~; seuls les ports nécessaires sont exposés vers l'hôte.

\begin{figure}[H]
  \centering
  \fbox{\parbox{0.88\textwidth}{\centering\vspace{3.2cm}
    \textit{[Diagramme d'architecture technique à services découplés
    --- à intégrer après export draw.io]}%
    \vspace{3.2cm}}}
  \caption{Architecture technique à services découplés de la plateforme EEC}
  \label{fig:diag-archi}
\end{figure}

%------------------------------------------------------------
\subsection{Modèle de Données}
\label{ssec:modele-bd}

Le schéma de données de la plateforme est un modèle relationnel-spatial
conforme aux normes PostGIS. Il est organisé en trois groupes d'entités
directement dérivés du domaine institutionnel de l'EEC et des exigences
fonctionnelles identifiées à la section~\ref{ssec:exigences}.

\noindent\textbf{Entités géographiques hiérarchiques.}
La structure territoriale de l'EEC est modélisée par quatre tables liées
en cascade~: \texttt{RegionSynodale}, qui stocke le polygone géospatial
de chacune des 22~régions (type \texttt{MultiPolygon}, SRID~4326)~;
\texttt{District}, rattaché par clé étrangère à une région~; \texttt{Paroisse},
rattachée à un district, avec un champ \texttt{position} de type
\texttt{Point} (SRID~4326) pour les coordonnées GPS, un indicateur
d'activité et les informations de contact~; \texttt{OEuvre}, rattachée à
une paroisse, avec un champ de type d'établissement (école, centre médical,
exploitation agropastorale, immeuble) et un champ de position facultatif.

\noindent\textbf{Ressources humaines.}
L'entité \texttt{Ouvrier} modélise le personnel ecclésiastique~: pasteur,
évangéliste ou diacre permanent, avec leur fonction et leur affectation à
une paroisse. L'entité \texttt{StatistiqueAnnuelle} est associée à une
paroisse et à une année~; elle enregistre les indicateurs démographiques
et financiers annuels.

\noindent\textbf{Gestion de la plateforme.}
L'entité \texttt{CompteUtilisateur} étend le modèle d'utilisateur standard~;
elle porte le rôle, les clés étrangères de périmètre synodal (région,
district ou paroisse), un champ de permissions personnalisées et un
indicateur de première connexion. L'entité \texttt{LogActivite} enregistre
de façon immuable chaque opération d'écriture~: auteur, type d'action,
modèle cible, identifiant de l'objet et horodatage UTC. Les entités visiteur
(\texttt{ProfilVisiteur}, \texttt{ParoisseEnregistree},
\texttt{ItinerairePersonnel}) complètent le schéma pour le module public
authentifié.

\begin{figure}[H]
  \centering
  \fbox{\parbox{0.88\textwidth}{\centering\vspace{3.2cm}
    \textit{[Schéma relationnel-spatial PostGIS
    --- à intégrer après génération draw.io]}%
    \vspace{3.2cm}}}
  \caption{Schéma de la base de données relationnelle-spatiale PostGIS}
  \label{fig:diag-bd}
\end{figure}

%------------------------------------------------------------
\subsection{Diagrammes de Séquence Système}
\label{ssec:seq-systeme}

Les diagrammes de séquence système modélisent les interactions entre les
acteurs externes et le système considéré comme une boîte noire~: ils montrent
les messages échangés et leur ordre chronologique, sans révéler les composants
internes. Deux scénarios représentatifs sont retenus.

\subsubsection{Séquence Système~1~--- Authentification et Accès au Tableau de Bord}

Ce diagramme (Figure~\ref{fig:seq-sys1}) représente le flux nominal de
connexion d'un administrateur~: saisie des identifiants, vérification par le
système, ouverture de session et redirection vers le tableau de bord adapté
au rôle. Il inclut le scénario alternatif de première connexion qui déclenche
obligatoirement le changement de mot de passe avant l'accès aux fonctions
d'administration.

\begin{figure}[H]
  \centering
  \fbox{\parbox{0.88\textwidth}{\centering\vspace{3.5cm}
    \textit{[Diagramme de séquence système 1
    --- Authentification et accès au tableau de bord
    --- à intégrer après génération draw.io]}%
    \vspace{3.5cm}}}
  \caption{Diagramme de séquence système --- Authentification et accès au tableau de bord}
  \label{fig:seq-sys1}
\end{figure}

\subsubsection{Séquence Système~2~--- Consultation Cartographique avec Filtrage}

Ce diagramme (Figure~\ref{fig:seq-sys2}) représente le flux de consultation
de la carte interactive par un utilisateur~: chargement de la carte, affichage
des couches institutionnelles, application d'un filtre hiérarchique par région
et district, mise à jour dynamique des marqueurs et consultation de la fiche
d'une entité sélectionnée. Il illustre la coordination entre l'interface
cartographique et le serveur de données géographiques.

\begin{figure}[H]
  \centering
  \fbox{\parbox{0.88\textwidth}{\centering\vspace{3.5cm}
    \textit{[Diagramme de séquence système 2
    --- Consultation cartographique avec filtrage hiérarchique
    --- à intégrer après génération draw.io]}%
    \vspace{3.5cm}}}
  \caption{Diagramme de séquence système --- Consultation cartographique avec filtrage}
  \label{fig:seq-sys2}
\end{figure}

%------------------------------------------------------------
\subsection{Diagrammes de Séquence Technique}
\label{ssec:seq-technique}

Les diagrammes de séquence technique descendent au niveau des composants
internes du système~: ils montrent les appels entre le frontend, le backend,
la base de données et le serveur cartographique, avec les noms de méthodes
et les flux de données réels.

\subsubsection{Séquence Technique~1~--- Gestion de Session Django}

Ce diagramme (Figure~\ref{fig:seq-tech1}) détaille le traitement interne
d'une requête d'authentification~: la validation des identifiants par le
backend, la création d'une entrée de session dans Redis, la génération du
cookie sécurisé \textit{HttpOnly} et sa transmission au navigateur, ainsi
que la vérification du jeton CSRF sur les requêtes modifiantes suivantes.
Il illustre la chaîne de sécurité complète mise en œuvre entre le frontend
Next.js et le backend Django.

\begin{figure}[H]
  \centering
  \fbox{\parbox{0.88\textwidth}{\centering\vspace{3.5cm}
    \textit{[Diagramme de séquence technique 1
    --- Authentification Django avec gestion de session et CSRF
    --- à intégrer après génération draw.io]}%
    \vspace{3.5cm}}}
  \caption{Diagramme de séquence technique --- Authentification Django et gestion de session}
  \label{fig:seq-tech1}
\end{figure}

\subsubsection{Séquence Technique~2~--- Requête Cartographique Leaflet vers GeoServer}

Ce diagramme (Figure~\ref{fig:seq-tech2}) détaille la chaîne de traitement
d'une requête cartographique émise par Leaflet.js~: construction de la
requête WMS par le client, transmission à GeoServer, interrogation de
PostGIS pour les entités géographiques, application du style SLD, génération
de la tuile image PNG et rendu dans l'interface. Il illustre le rôle de
GeoServer comme couche intermédiaire entre la base de données spatiale et
le client cartographique.

\begin{figure}[H]
  \centering
  \fbox{\parbox{0.88\textwidth}{\centering\vspace{3.5cm}
    \textit{[Diagramme de séquence technique 2
    --- Requête WMS Leaflet → GeoServer → PostGIS
    --- à intégrer après génération draw.io]}%
    \vspace{3.5cm}}}
  \caption{Diagramme de séquence technique --- Requête cartographique WMS de Leaflet vers GeoServer}
  \label{fig:seq-tech2}
\end{figure}

%------------------------------------------------------------
\subsection{Diagramme de Classes Technique}
\label{ssec:classes-tech}

Le diagramme de classes technique (Figure~\ref{fig:diag-classes-tech})
représente le modèle d'implémentation de la plateforme~: il étend le
diagramme de classes métier (Figure~\ref{fig:diag-classes-metier}) avec
les détails d'implémentation propres à Django et à GeoDjango~:
types de données précis (\texttt{PointField}, \texttt{MultiPolygonField},
\texttt{JSONField}), visibilité des attributs, signatures des méthodes
(sérialiseurs DRF, méthodes de filtrage RBAC, helpers de calcul spatial).
Les dépendances entre les classes de modèles, les sérialiseurs et les vues
DRF sont représentées explicitement.

\begin{figure}[H]
  \centering
  \fbox{\parbox{0.88\textwidth}{\centering\vspace{3.2cm}
    \textit{[Diagramme de classes technique
    --- à intégrer après génération draw.io]}%
    \vspace{3.2cm}}}
  \caption{Diagramme de classes technique de l'implémentation Django/GeoDjango}
  \label{fig:diag-classes-tech}
\end{figure}

\medskip
\noindent\textbf{Bilan du chapitre.}
Ce chapitre a formalisé l'ensemble des besoins de la plateforme en
dix-huit exigences fonctionnelles et dix exigences non-fonctionnelles, puis
modélisé le système à travers six diagrammes d'analyse --- contexte,
packages, cas d'utilisation, classes métier et activité --- et cinq
diagrammes de conception --- architecture technique, schéma de données
PostGIS, deux diagrammes de séquence système, deux diagrammes de séquence
technique et diagramme de classes technique. Ces modèles constituent le
contrat de spécification sur lequel s'appuie l'implémentation détaillée au
\textbf{Chapitre~3}.

\clearpage

% ============================================================
%  CHAPITRE 3 — IMPLÉMENTATION ET RÉSULTATS
% ============================================================
\chapter{Implémentation et Résultats}
\label{chap:implementation}

\begin{chapplan}
  \item \textbf{Section~\ref{sec:outils} --- Outils et Technologies}~:
    justification des choix technologiques pour chaque composant de la
    pile logicielle~; tableau récapitulatif des outils, versions et rôles.
  \item \textbf{Section~\ref{sec:deploiement} --- Diagramme de Déploiement}~:
    distribution des cinq services sur les nœuds physiques et virtuels~;
    configuration réseau Docker Compose.
  \item \textbf{Section~\ref{sec:resultats} --- Résultats}~:
    présentation annotée des interfaces fonctionnelles~: carte interactive,
    authentification, tableaux de bord par rôle, import/export et journal
    d'audit~; confrontation aux exigences du Chapitre~2.
  \item \textbf{Section~\ref{sec:discussion} --- Discussion}~:
    vérification des objectifs, analyse critique des forces et des limites
    de la solution réalisée.
\end{chapplan}

\lettrine[lines=2, loversize=0.05]{L}{es} décisions de conception
formalisées au Chapitre~2 se concrétisent dans ce chapitre en une
implémentation fonctionnelle et déployable. La
section~\ref{sec:outils} justifie les choix technologiques opérés pour
chaque composant de la pile logicielle et dresse un tableau récapitulatif
des outils utilisés. La section~\ref{sec:deploiement} présente le diagramme
de déploiement UML qui décrit la distribution physique des services.
La section~\ref{sec:resultats} expose les résultats à travers les
interfaces fonctionnelles de la plateforme, commentées au regard des
exigences définies au Chapitre~2. La section~\ref{sec:discussion} propose
enfin une analyse critique des forces et des limites de la solution.

%------------------------------------------------------------
\section{Outils et Technologies}
\label{sec:outils}
%------------------------------------------------------------

\subsection{Justification des Choix Technologiques}
\label{ssec:justification}

Le choix de chaque technologie résulte d'une évaluation explicite des
alternatives disponibles au regard des contraintes identifiées~: contraintes
institutionnelles (déploiement sur infrastructure limitée, données
sensibles), contraintes fonctionnelles (requêtes spatiales, RBAC
hiérarchique, publication OGC) et contraintes de maintenabilité
(logiciels libres, documentation active, communauté de support).

\subsubsection{Django~5 et Django REST Framework}

Django~5~\cite{django2024} a été retenu comme framework backend principal
pour trois raisons déterminantes. En premier lieu, son extension spatiale
native GeoDjango offre une intégration directe avec PostGIS~: les requêtes
de distance (\texttt{Distance()}), les filtres spatiaux (\texttt{contains},
\texttt{intersects}) et les champs de géométrie (\texttt{PointField},
\texttt{MultiPolygonField}) sont disponibles via l'ORM sans SQL manuel.
En second lieu, le module d'authentification de Django~5 propose une gestion
de sessions robuste par cookies \textit{HttpOnly}~; ce mécanisme est
nativement protégé contre les attaques de type XSS, contrairement aux
jetons JWT transmis dans le stockage local du navigateur, qui présentaient
une vulnérabilité documentée dans la version antérieure de la plateforme.
En troisième lieu, l'écosystème Python disponible ---
\texttt{openpyxl} pour la lecture des fichiers Excel EEC,
\texttt{reportlab} pour la génération PDF, \texttt{django-cors-headers}
pour la politique d'origines croisées --- couvre l'ensemble des besoins
fonctionnels sans dépendances tierces lourdes. Django REST
Framework~\cite{drf2024} complète Django en exposant les ressources
institutionnelles sous forme d'API REST JSON consommée par le frontend.

\subsubsection{PostgreSQL~16 et PostGIS~3.4}

PostgreSQL~16 associé à PostGIS~3.4~\cite{postgis2023} constitue le
standard de facto pour les bases de données spatiales open source~: il est
utilisé dans des projets cartographiques institutionnels à l'échelle nationale
et internationale, dispose d'une indexation spatiale performante
(\textit{R-tree GiST}) et garantit la conformité aux spécifications OGC
Simple Features. La fonction \texttt{ST\_Distance()} native permet
d'identifier la paroisse la plus proche d'une position GPS en une seule
requête SQL, sans traitement applicatif. Aucun SGBD alternatif ne combine
ces fonctionnalités spatiales et une intégration aussi mature avec GeoDjango.

\subsubsection{GeoServer~2.26}

GeoServer~2.26~\cite{geoserver2024} a été retenu pour la publication des
couches cartographiques institutionnelles au regard de sa conformité native
aux protocoles OGC WMS~1.3.0 et WFS~2.0.0, qui garantissent
l'interopérabilité avec tout client cartographique standard. La
configuration des styles visuels via des fichiers SLD permet de définir les
couleurs institutionnelles de l'EEC (vert synodal, fond clair) directement
sur le serveur, sans traitement côté client. L'alternative d'une
publication de fichiers GeoJSON statiques, utilisée dans la version
antérieure, était incompatible avec la mise à jour dynamique des couches
et ne répondait pas aux exigences ENF05 d'interopérabilité OGC.

\subsubsection{Next.js~14}

Next.js~14~\cite{nextjs2024} avec l'App Router a été retenu pour le frontend
pour sa capacité à combiner rendu côté serveur (\textit{SSR}) et rendu côté
client (\textit{CSR}) dans la même application. Le \textit{SSR} est
utilisé pour les pages d'administration~: les données sensibles ne sont
jamais envoyées au navigateur avant vérification de la session, ce qui
réduit la surface d'attaque. Le rendu côté client est utilisé pour la carte
interactive, qui requiert des interactions dynamiques non compatibles avec
le \textit{SSR}. L'App Router de Next.js~14 permet de définir des
\textit{layouts} partagés par rôle d'administration (SUPER/REGION,
DISTRICT, PAROISSE), ce qui simplifie l'implémentation du RBAC côté
frontend.

\subsubsection{Leaflet.js et MarkerCluster}

Leaflet.js~\cite{leaflet2024} a été retenu pour la bibliothèque
cartographique cliente pour sa légèreté (environ 42~Ko minifié), sa licence
libre et l'étendue de son écosystème de plugins. Le plugin
\texttt{Leaflet.MarkerCluster} implémente le regroupement automatique des
565~marqueurs de paroisses en indicateurs numériques~: sans ce mécanisme,
la densité des marqueurs rendrait la carte illisible aux niveaux de zoom
nationaux. La prise en charge native des couches \textit{TileLayer.WMS}
permet de consommer directement les couches GeoServer sans adaptation.

\subsubsection{Redis~7 et Docker Compose}

Redis~7~\cite{redis2024} assure deux fonctions complémentaires~:
le stockage des sessions Django côté serveur (évitant leur persistance
en base de données relationnelle) et la mise en cache des réponses API
fréquentes. Docker Compose~\cite{docker2024} orchestre l'ensemble des
cinq services dans une configuration déclarative reproductible sur toute
machine hôte disposant de Docker~: c'est la réponse directe à l'exigence
ENF06 de portabilité et à la limite de déploiement identifiée dans la
plateforme GeoEEC~v1.0.

%-----
\subsection{Récapitulatif des Outils Utilisés}
\label{ssec:outils-table}

Le Tableau~\ref{tab:outils} dresse la liste complète des outils et
technologies utilisés dans le projet, avec leur version exacte et leur
rôle précis.

\begin{table}[H]
\centering
\caption{Récapitulatif des outils et technologies de la plateforme EEC}
\label{tab:outils}
\setlength{\tabcolsep}{5pt}
\renewcommand{\arraystretch}{1.22}
\resizebox{\linewidth}{!}{%
\begin{tabular}{%
  >{\bfseries}m{3.6cm}
  >{\centering\arraybackslash}m{1.6cm}
  m{9.0cm}}
\toprule
\rowcolor{eecgreen}
\multicolumn{1}{>{\bfseries\color{white}}m{3.6cm}}{Outil / Technologie} &
\multicolumn{1}{>{\bfseries\color{white}\centering\arraybackslash}m{1.6cm}}{Version} &
\multicolumn{1}{>{\bfseries\color{white}}m{9.0cm}}{Rôle dans le projet} \\
\midrule
\rowcolor{eeclightgreen!60}
Python           & 3.12   & Langage principal du backend \\
Django           & 5.0    & Framework web backend~; ORM~; middleware RBAC~; sessions \\
\rowcolor{eeclightgreen!60}
GeoDjango        & 5.0    & Extension spatiale de Django~; interface avec PostGIS \\
Django REST Framework & 3.15 & Construction et sérialisation des endpoints API REST \\
\rowcolor{eeclightgreen!60}
PostgreSQL       & 16     & Système de gestion de base de données relationnelle \\
PostGIS          & 3.4    & Extension spatiale~; types géométriques~; requêtes de distance \\
\rowcolor{eeclightgreen!60}
GeoServer        & 2.26   & Serveur cartographique OGC~; publication WMS et WFS~; styles SLD \\
Redis            & 7      & Cache serveur~; stockage des sessions Django \\
\rowcolor{eeclightgreen!60}
Next.js          & 14     & Framework frontend React~; App Router~; SSR et CSR \\
TypeScript       & 5.x    & Typage statique du code frontend \\
\rowcolor{eeclightgreen!60}
Tailwind CSS     & 3.x    & Framework de mise en forme CSS utilitaire \\
Leaflet.js       & 1.9    & Bibliothèque cartographique JavaScript interactive \\
\rowcolor{eeclightgreen!60}
Leaflet.MarkerCluster & 1.5 & Regroupement automatique des marqueurs sur la carte \\
openpyxl         & 3.1    & Lecture et écriture des fichiers Excel (.xlsx) \\
\rowcolor{eeclightgreen!60}
reportlab        & 4.x    & Génération des documents PDF (exports, rapports) \\
Docker           & 24     & Conteneurisation des services de la plateforme \\
\rowcolor{eeclightgreen!60}
Docker Compose   & 2.x    & Orchestration des cinq services en environnement multi-conteneurs \\
draw.io / diagrams.net & --- & Conception des diagrammes UML du rapport \\
\bottomrule
\end{tabular}}
\end{table}

%------------------------------------------------------------
\section{Diagramme de Déploiement}
\label{sec:deploiement}
%------------------------------------------------------------

Le diagramme de déploiement (Figure~\ref{fig:diag-deploiement}) représente
la distribution physique et virtuelle des composants logiciels de la
plateforme sur les nœuds d'exécution. Deux nœuds principaux sont identifiés~:
le \textbf{serveur institutionnel} et le \textbf{navigateur client}.

Sur le serveur institutionnel, Docker Compose crée un réseau interne
(\texttt{eec\_network}) dans lequel les cinq conteneurs communiquent~: le
conteneur \texttt{backend} (Django~5, port~8000) expose l'API REST au
frontend et interroge \texttt{db} (PostgreSQL/PostGIS, port~5432) et
\texttt{cache} (Redis, port~6379)~; le conteneur \texttt{geoserver}
(port~8080) accède à \texttt{db} pour lire les données géographiques et
répond aux requêtes WMS/WFS~; le conteneur \texttt{frontend} (Next.js~14,
port~3000) sert l'interface web. Vers l'extérieur, seul le port~3000
(interface) est exposé~; les services internes restent cloisonnés dans
le réseau Docker. Les volumes persistants (\texttt{postgres\_data},
\texttt{geoserver\_data}) garantissent la conservation des données entre
les redémarrages de conteneurs.

Sur le navigateur client, Leaflet.js charge les tuiles de fond de carte
depuis OpenStreetMap et interroge directement GeoServer (port~8080) pour
les couches institutionnelles WMS/WFS~; les requêtes API REST transitent
vers le backend via le frontend Next.js.

\begin{figure}[H]
  \centering
  \fbox{\parbox{0.88\textwidth}{\centering\vspace{3.2cm}
    \textit{[Diagramme de déploiement UML
    --- à intégrer après génération draw.io]}%
    \vspace{3.2cm}}}
  \caption{Diagramme de déploiement UML --- Distribution des services Docker Compose}
  \label{fig:diag-deploiement}
\end{figure}

%------------------------------------------------------------
\section{Résultats}
\label{sec:resultats}
%------------------------------------------------------------

Cette section présente les interfaces fonctionnelles de la plateforme
implémentée, commentées au regard des exigences définies au
Chapitre~2. Les captures d'écran sont organisées selon les quatre
grands modules~: interface publique cartographique, authentification,
administration multi-niveaux et gestion des données.

%-----
\subsection{Interface Publique Cartographique}
\label{ssec:res-carte}

L'interface publique est accessible sans authentification depuis la racine
de la plateforme. Elle répond directement aux exigences EF03, EF04, EF05
et EF16 relatives à la cartographie interactive.

La Figure~\ref{fig:res-carte-publique} illustre la carte au niveau de zoom
national~: les 553~paroisses importées sont regroupées en indicateurs de
cluster numériques selon leur densité géographique. Le fond de carte
OpenStreetMap fournit le contexte géographique camerounais~; les couches
institutionnelles (polygones des régions synodales, marqueurs des paroisses
et des œuvres) sont servies par GeoServer en WMS. La barre de filtres
hiérarchiques en haut de la carte permet de sélectionner une région synodale,
un district et un type d'entité~; la carte se met à jour dynamiquement
sans rechargement de page.

\begin{figure}[H]
  \centering
  % ===== IMAGE 1 : Carte publique avec clustering =====
  \fbox{\parbox{0.88\textwidth}{\centering\vspace{3.8cm}
    \textit{[Capture d'écran 1 --- Carte interactive publique avec clustering de marqueurs]}%
    \vspace{3.8cm}}}
  \caption{Interface publique --- Carte interactive avec regroupement automatique des marqueurs}
  \label{fig:res-carte-publique}
\end{figure}

La Figure~\ref{fig:res-popup} montre le panneau de détail qui s'affiche lors
d'un clic sur un marqueur individuel~: nom de la paroisse, district et région
d'appartenance, coordonnées GPS, contacts et statut d'activité. Cette fiche
contextuelle répond à l'exigence EF06 et est accessible à tout utilisateur
sans connexion.

\begin{figure}[H]
  \centering
  % ===== IMAGE 2 : Popup de détail d'une paroisse =====
  \fbox{\parbox{0.88\textwidth}{\centering\vspace{3.8cm}
    \textit{[Capture d'écran 2 --- Popup de détail d'une paroisse sur la carte]}%
    \vspace{3.8cm}}}
  \caption{Interface publique --- Popup de détail d'une paroisse sélectionnée}
  \label{fig:res-popup}
\end{figure}

%-----
\subsection{Interface d'Authentification}
\label{ssec:res-auth}

La Figure~\ref{fig:res-login} présente la page de connexion de l'espace
d'administration (\texttt{/admin/login}). Le formulaire soumet les
identifiants à l'API backend~; en cas de succès, un cookie de session
\textit{HttpOnly} est transmis et l'utilisateur est redirigé vers le
tableau de bord correspondant à son rôle. En cas d'identifiants incorrects,
un message d'erreur est affiché sans révéler si l'identifiant ou le mot de
passe est la source de l'erreur, conformément aux bonnes pratiques de
sécurité recommandées par l'OWASP~\cite{owasp2021}. Pour un compte dont
l'indicateur de première connexion est actif, la redirection s'effectue
automatiquement vers le formulaire de changement de mot de passe.

\begin{figure}[H]
  \centering
  % ===== IMAGE 3 : Page de connexion administrateur =====
  \fbox{\parbox{0.88\textwidth}{\centering\vspace{3.8cm}
    \textit{[Capture d'écran 3 --- Page de connexion de l'espace administrateur]}%
    \vspace{3.8cm}}}
  \caption{Interface d'authentification --- Page de connexion sécurisée}
  \label{fig:res-login}
\end{figure}

%-----
\subsection{Interface d'Administration Multi-Niveaux}
\label{ssec:res-admin}

L'interface d'administration adapte son contenu et ses fonctionnalités au
rôle de l'utilisateur connecté, conformément aux exigences EF02 et EF07 à
EF09. Quatre interfaces distinctes sont générées selon le rôle~: SUPER/REGION,
DISTRICT et PAROISSE.

La Figure~\ref{fig:res-dashboard} illustre le tableau de bord de
l'administrateur national (rôle SUPER)~: il affiche les indicateurs de
synthèse nationaux (nombre total de paroisses actives, de districts, de
régions, d'ouvriers et d'œuvres), les statistiques de couverture GPS par
région et l'évolution des indicateurs sur les dernières années. Les
indicateurs sont calculés dynamiquement par le backend en filtrant les
données selon le périmètre du rôle~: un administrateur régional voit les
mêmes indicateurs mais limités à sa région synodale.

\begin{figure}[H]
  \centering
  % ===== IMAGE 4 : Tableau de bord administrateur national =====
  \fbox{\parbox{0.88\textwidth}{\centering\vspace{3.8cm}
    \textit{[Capture d'écran 4 --- Tableau de bord de l'administrateur national]}%
    \vspace{3.8cm}}}
  \caption{Interface d'administration --- Tableau de bord de l'administrateur national}
  \label{fig:res-dashboard}
\end{figure}

La Figure~\ref{fig:res-paroisses} présente l'interface de gestion des
paroisses~: liste paginée filtrée par région et district, formulaire de
création et de modification avec validation des coordonnées GPS en temps
réel, et action de suppression avec confirmation. Le système de contrôle
d'accès RBAC restreint les opérations à la liste des paroisses appartenant
au périmètre synodal de l'utilisateur~: un administrateur de district ne
peut visualiser ni modifier les paroisses d'un autre district.

\begin{figure}[H]
  \centering
  % ===== IMAGE 5 : Interface de gestion des paroisses =====
  \fbox{\parbox{0.88\textwidth}{\centering\vspace{3.8cm}
    \textit{[Capture d'écran 5 --- Interface de gestion des paroisses (liste et formulaire)]}%
    \vspace{3.8cm}}}
  \caption{Interface d'administration --- Gestion des paroisses avec contrôle RBAC}
  \label{fig:res-paroisses}
\end{figure}

%-----
\subsection{Module d'Import et d'Export}
\label{ssec:res-import}

Le module d'import (Figure~\ref{fig:res-import}) répond à l'exigence EF10~:
il permet de charger les fichiers Excel fournis par les services de l'EEC
directement dans la base de données, après une validation ligne par ligne.
Le rapport d'import affiché après traitement détaille le nombre de lignes
insérées avec succès, le nombre de lignes rejetées et, pour chaque
rejet, la raison précise (coordonnées GPS hors territoire, district
inexistant, doublon détecté). Cette traçabilité de l'import répond
directement au problème de qualité des données identifié lors de l'audit
préliminaire~: 127~paroisses sources sans coordonnées GPS valides ont ainsi
été identifiées et signalées lors de la première importation.

\begin{figure}[H]
  \centering
  % ===== IMAGE 6 : Module d'import Excel et rapport de résultats =====
  \fbox{\parbox{0.88\textwidth}{\centering\vspace{3.8cm}
    \textit{[Capture d'écran 6 --- Module d'import Excel avec rapport détaillé de résultats]}%
    \vspace{3.8cm}}}
  \caption{Module d'import --- Chargement de données Excel et rapport de validation}
  \label{fig:res-import}
\end{figure}

%-----
\subsection{Journal d'Audit}
\label{ssec:res-audit}

Le journal d'audit (Figure~\ref{fig:res-audit}) est accessible
exclusivement à l'administrateur national. Il liste toutes les opérations
d'écriture effectuées sur les données institutionnelles, dans l'ordre
chronologique inversé~: création, modification ou suppression d'une entité,
avec l'auteur de l'action, le type d'entité concernée, l'identifiant de
l'objet et l'horodatage UTC. Des filtres par type d'action, type d'entité
et période permettent de cibler une recherche. En cliquant sur une entrée,
l'administrateur accède au détail de la modification~: valeurs avant et
après la mise à jour. Ce module répond aux exigences EF14 et EF15 relatives
à la traçabilité, et constitue l'un des apports distinctifs par rapport
à la plateforme GeoEEC~v1.0 qui ne disposait d'aucun mécanisme d'audit.

\begin{figure}[H]
  \centering
  % ===== IMAGE 7 : Journal d'audit =====
  \fbox{\parbox{0.88\textwidth}{\centering\vspace{3.8cm}
    \textit{[Capture d'écran 7 --- Journal d'audit des modifications institutionnelles]}%
    \vspace{3.8cm}}}
  \caption{Journal d'audit --- Historique des opérations d'écriture sur les données}
  \label{fig:res-audit}
\end{figure}

%------------------------------------------------------------
\section{Discussion}
\label{sec:discussion}
%------------------------------------------------------------

\subsection{Adéquation aux Objectifs}
\label{ssec:disc-objectifs}

Les six objectifs formulés en Introduction Générale sont confrontés aux
résultats obtenus. L'objectif~1 (modélisation du schéma de données
hiérarchique) est atteint~: le schéma PostGIS implémenté couvre les
22~régions synodales, 137~districts, 553~paroisses importées, 311~œuvres
et 685~ouvriers, avec toutes les contraintes d'intégrité référentielle et
les champs de géolocalisation. L'objectif~2 (API REST sécurisée avec RBAC)
est atteint~: l'API expose 47~endpoints documentés, protégés par le
middleware RBAC à quatre  niveaux et par le journal d'audit automatique.
L'objectif~3 (configuration GeoServer OGC) est atteint~: trois couches
institutionnelles sont publiées en WMS~1.3.0 et WFS~2.0.0 avec styles SLD.
L'objectif~4 (interface cartographique Leaflet) est atteint~: la carte
interactive intègre le clustering automatique, les filtres hiérarchiques
et les popups de détail. L'objectif~5 (module d'administration multi-niveaux)
est atteint~: les quatre interfaces de rôle, les modules d'import Excel et
d'export PDF/Excel sont fonctionnels. L'objectif~6 (conteneurisation Docker)
est atteint~: la plateforme est entièrement déployable avec une commande
unique sur tout serveur disposant de Docker.

%-----
\subsection{Forces de la Solution}
\label{ssec:disc-forces}

\textbf{Sécurité améliorée par rapport à la version antérieure.}
Le remplacement de l'authentification JWT par des sessions Django avec
cookies \textit{HttpOnly} et \textit{SameSite=Strict} élimine le vecteur
d'attaque XSS qui permettait la capture des jetons transmis dans le stockage
local du navigateur. La protection CSRF systématique sur toutes les
requêtes modifiantes renforce ce dispositif, conformément aux recommandations
de l'OWASP~\cite{owasp2021}.

\textbf{Interopérabilité cartographique.}
La publication des couches géographiques via GeoServer en WMS~1.3.0 et
WFS~2.0.0 rend les données de l'EEC consommables par tout client
cartographique compatible OGC~: QGIS, OpenLayers, ou un futur portail
institutionnel camerounais. Aucune solution concurrente analysée au
Chapitre~1 ne proposait cette combinaison.

\textbf{RBAC aligné sur la hiérarchie réelle de l'EEC.}
Le modèle de contrôle d'accès à quatre niveaux --- national, régional,
district et paroissial --- correspond exactement à la structure synodale
de l'institution~: un administrateur de district peut gérer les paroisses
de son périmètre sans accès aux données des districts voisins. Ce niveau
de granularité était absent de la version antérieure.

\textbf{Traçabilité complète et automatique.}
Tout enregistrement de modification, de création ou de suppression
déclenche automatiquement l'écriture d'une entrée dans le journal d'audit
via les signaux Django, sans intervention du développeur dans chaque vue.
Cette automatisation garantit l'exhaustivité de l'historique, impossible
à assurer manuellement.

\textbf{Déploiement reproductible.}
La configuration Docker Compose en cinq services garantit que l'environnement
d'exécution est identique en développement et en production~: plus aucune
divergence de versions de bibliothèques ou de configuration système ne peut
causer des défaillances à l'installation.

Ces cinq forces confirment la pertinence des décisions architecturales
formalisées au Chapitre~2. Elles positionnent la plateforme de façon
distincte par rapport aux solutions analysées dans l'état de l'art, en
particulier sur les axes de sécurité, d'interopérabilité OGC et de
RBAC institutionnel. La section suivante expose les limites structurelles
qui tempèrent la portée de ces résultats.

%------------------------------------------------------------
\section{Limites et Contraintes}
\label{sec:limites}
%------------------------------------------------------------

Tout travail d'ingénierie comporte des contraintes qui en délimitent la
portée. Une évaluation honnête de ces limites est indispensable pour
apprécier correctement la valeur de la solution et orienter les
développements futurs.

\textbf{Couverture GPS incomplète.}
L'audit des données sources a révélé que 127~paroisses sur 565, soit
22~\% de l'effectif total, ne disposent d'aucune coordonnée GPS valide
dans les fichiers Excel fournis par l'EEC. Ces paroisses sont importées
en base de données mais ne s'affichent pas sur la carte interactive.
La plateforme ne peut donc pas prétendre à une couverture cartographique
complète du réseau ecclésiastique~; son utilité pour les décisions de
planification territoriale nationale reste limitée tant que ce taux
n'est pas réduit par une campagne de collecte terrain.

\textbf{Distance à vol d'oiseau.}
La fonction de calcul d'itinéraire repose sur la formule de Haversine,
qui mesure la distance géodésique en ligne droite entre deux points.
Cette mesure sous-estime systématiquement la durée réelle de trajet
sur un réseau routier camerounais caractérisé par des voies sinueuses
et un relief accidenté~; elle ne tient pas compte des barrières
physiques (fleuves, massifs montagneux, absence de pont).
La durée estimée à 60~km/h de moyenne est donc approximative et ne
peut servir de base à des décisions opérationnelles précises.

\textbf{Absence d'application mobile.}
La plateforme est conçue pour des écrans desktop et tablette
(résolution minimale 1\,024~$\times$~768~px, ENF09). Elle n'est pas
adaptée aux smartphones, qui constituent pourtant le principal outil
numérique des administrateurs paroissiaux en zone rurale camerounaise.
En l'absence d'application mobile dédiée, la collecte terrain des
coordonnées GPS des 127~paroisses non géolocalisées reste sans
outillage numérique adapté.

\textbf{Authentification à facteur unique.}
L'accès à l'espace d'administration repose sur un couple
identifiant/mot de passe sans mécanisme de double authentification
(TOTP ou code temporaire par messagerie). Pour des comptes de niveau
national disposant d'un accès complet à l'ensemble des données de
l'EEC, cette limite expose la plateforme aux risques de compromission
par capture ou réutilisation de mot de passe, en l'absence de second
facteur de vérification.

\textbf{Déploiement mono-instance.}
L'architecture Docker Compose actuelle déploie chaque service en
instance unique sur un seul serveur physique. Elle ne prévoit pas de
mécanisme de réplication ou de bascule automatique en cas de défaillance
matérielle~; la disponibilité de la plateforme est donc directement
dépendante de la fiabilité du serveur institutionnel hôte.

\medskip
\noindent\textbf{Bilan du chapitre.}
Ce chapitre a présenté et justifié les choix technologiques de la
plateforme à travers quatre tableaux comparatifs, exposé les résultats
à travers sept interfaces fonctionnelles commentées, analysé les cinq
forces distinctives de la solution dans la section Discussion, puis
identifié cinq limites structurelles dans cette section. La
\textbf{Conclusion Générale} synthétise l'ensemble de la démarche
et formule les perspectives d'amélioration directement dérivées
de ces limites.

\clearpage

% ============================================================
%  CONCLUSION GÉNÉRALE
% ============================================================
\chapter*{Conclusion Générale}
\addcontentsline{toc}{chapter}{Conclusion Générale}
\markboth{Conclusion Générale}{Conclusion Générale}

\noindent
L'Église Évangélique du Cameroun gérait jusqu'alors ses 565~paroisses,
311~œuvres sociales et 685~ouvriers ecclésiastiques répartis sur
22~régions synodales et 137~districts sans aucun outil numérique
centralisé~: pas de visualisation géographique, pas de contrôle d'accès
formalisé selon les périmètres synodaux, pas de journal d'audit des
modifications, pas de publication cartographique interopérable. Ce travail
a abordé ce problème par la conception et le développement d'une
plateforme web cartographique institutionnelle, en s'appuyant sur
l'hypothèse qu'une architecture à services découplés, entièrement
conteneurisée, combinant un serveur cartographique OGC et un contrôle
d'accès hiérarchique aligné sur la structure synodale, permettrait d'y
répondre de façon complète et déployable dans un contexte institutionnel
à ressources limitées.

Cette hypothèse est validée par les résultats obtenus~: la plateforme
implémentée déploie en une seule commande cinq services
--- Django~5~+~GeoDjango, PostgreSQL~+~PostGIS~3.4, GeoServer~2.26,
Next.js~14 et Redis~7~--- et offre une carte interactive avec clustering
de 553~paroisses, des filtres hiérarchiques par région et district, un
tableau de bord adaptatif par rôle, un module d'import/export des données
institutionnelles, et un journal d'audit automatique et exhaustif de toutes
les opérations d'écriture. L'ensemble des six objectifs formulés en
Introduction Générale a été atteint. Par rapport à la plateforme
GeoEEC~v1.0, la solution apporte une valeur ajoutée précise et documentée
sur cinq axes~: sécurité de l'authentification par sessions \textit{HttpOnly},
RBAC aligné sur les quatre niveaux de la hiérarchie synodale réelle,
publication cartographique conforme aux standards OGC WMS et WFS, traçabilité
automatique par journal d'audit, et reproductibilité de l'environnement
d'exécution par conteneurisation complète.

\medskip
\noindent\textbf{Limites.}
Cinq limites structurelles, analysées en détail dans la
section~\ref{sec:limites} du Chapitre~3, tempèrent la portée de ces
résultats. La couverture cartographique reste incomplète~: 127~paroisses
(22~\%) ne disposent d'aucune coordonnée GPS dans les fichiers sources de
l'EEC et ne s'affichent pas sur la carte, ce qui restreint directement
l'utilité de la plateforme pour les décisions de planification territoriale.
Le calcul d'itinéraire repose sur la distance à vol d'oiseau et non sur
le réseau routier réel, produisant des durées de trajet sous-estimées dans
les zones à relief accidenté. L'absence d'interface mobile prive les
administrateurs paroissiaux d'un outil de saisie GPS sur le terrain.
L'authentification à facteur unique expose les comptes nationaux à des
risques de compromission par mot de passe seul. Enfin, le déploiement
mono-instance ne prévoit pas de mécanisme de haute disponibilité en cas
de défaillance du serveur hôte.

\medskip
\noindent\textbf{Perspectives.}
Ces limites définissent directement les axes de développement prioritaires.
La première perspective est le développement d'une application mobile
légère (React Native ou Progressive Web App) permettant aux administrateurs
paroissiaux de saisir ou de corriger les coordonnées GPS directement sur le
terrain~: cette seule fonctionnalité permettrait de résorber le déficit de
géolocalisation des 127~paroisses non couvertes et de faire passer le taux
de couverture cartographique de 78~\% à 100~\%. La deuxième perspective est
l'intégration d'un moteur de calcul d'itinéraires routiers~--- via l'API
OSRM (\textit{Open Source Routing Machine}) ou une instance locale de
GraphHopper~--- pour remplacer la distance à vol d'oiseau par des temps de
trajet réalistes sur les routes camerounaises. La troisième perspective est
l'activation de la double authentification par TOTP (\textit{Time-based
One-Time Password}) pour les comptes d'administrateurs nationaux et
régionaux, en s'appuyant sur la bibliothèque \texttt{django-otp}
déjà compatible avec l'infrastructure Django. Enfin, la réintégration
du module de questions-réponses en langage naturel développé dans
GeoEEC~v1.0 --- enrichi des nouvelles entités et des statistiques annuelles
--- ouvrirait la plateforme à des usages analytiques avancés permettant
aux responsables synodaux d'interroger les données institutionnelles sans
compétences techniques.

\noindent\rule{\linewidth}{0.8pt}
 dans complet.tex,
%            juste après le contenu de la Section 2.2 existante.
% ============================================================

% ============================================================
%  SUITE — SECTION 2.2 CONCEPTION DU SYSTÈME
% ============================================================

\subsection{Architecture Technique}
\label{ssec:archi}

L'architecture retenue pour la plateforme EEC~Géolocalisation est une
\textbf{architecture à services découplés}, dans laquelle chaque composant
fonctionnel est isolé dans un conteneur indépendant et communique avec les
autres via des interfaces standardisées. Ce choix s'oppose à une architecture
monolithique dans laquelle toutes les fonctions seraient regroupées dans un
seul processus~: la séparation des responsabilités facilite la maintenance,
permet la mise à jour indépendante de chaque couche et renforce la sécurité en
cloisonnant les surfaces d'attaque.

La plateforme est décomposée en cinq services orchestrés par Docker~Compose,
représentés dans la Figure~\ref{fig:diag-archi}~:

\begin{itemize}[itemsep=3pt, topsep=5pt, leftmargin=1.5cm]
  \item \textbf{Backend} --- serveur Django~5 avec GeoDjango et
    Django REST Framework~: expose l'intégralité de la logique métier
    sous forme d'API REST, applique le contrôle d'accès basé sur les rôles
    et gère les sessions d'authentification par cookies sécurisés
    (\textit{HttpOnly}, \textit{SameSite=Strict})~;
  \item \textbf{Base de données} --- PostgreSQL~16 avec l'extension
    PostGIS~3.4~: assure la persistance des données institutionnelles et
    des géométries spatiales (points, polygones)~; les requêtes de
    proximité exploitent les fonctions de calcul de distance natives~;
  \item \textbf{Serveur cartographique} --- GeoServer~2.26~: publie les
    couches géographiques de l'EEC selon les standards OGC WMS~1.3.0 et
    WFS~2.0.0~; les styles visuels sont définis via des fichiers SLD aux
    couleurs institutionnelles~;
  \item \textbf{Cache} --- Redis~7~: stocke les sessions Django côté
    serveur et met en cache les réponses fréquentes de l'API pour réduire
    la charge sur la base de données~;
  \item \textbf{Frontend} --- Next.js~14 avec l'App Router~: sert
    l'interface publique (carte Leaflet, clustering, filtres) et
    l'interface d'administration multi-niveaux~; les pages critiques sont
    rendues côté serveur (\textit{SSR}) pour garantir la sécurité des
    données sensibles.
\end{itemize}

Les échanges entre le frontend et le backend transitent exclusivement par
l'API REST en JSON. Le frontend Leaflet interroge GeoServer directement
pour les couches cartographiques, sans passer par le backend Django~: cette
séparation permet d'exploiter les protocoles OGC standard sans surcharger
le serveur applicatif. L'ensemble des services partage un réseau Docker
interne isolé~; seuls les ports nécessaires sont exposés vers l'hôte.

\begin{figure}[H]
  \centering
  \fbox{\parbox{0.88\textwidth}{\centering\vspace{3.2cm}
    \textit{[Diagramme d'architecture technique à services découplés
    --- à intégrer après export draw.io]}%
    \vspace{3.2cm}}}
  \caption{Architecture technique à services découplés de la plateforme EEC}
  \label{fig:diag-archi}
\end{figure}

%------------------------------------------------------------
\subsection{Modèle de Données}
\label{ssec:modele-bd}

Le schéma de données de la plateforme est un modèle relationnel-spatial
conforme aux normes PostGIS. Il est organisé en trois groupes d'entités
directement dérivés du domaine institutionnel de l'EEC et des exigences
fonctionnelles identifiées à la section~\ref{ssec:exigences}.

\noindent\textbf{Entités géographiques hiérarchiques.}
La structure territoriale de l'EEC est modélisée par quatre tables liées
en cascade~: \texttt{RegionSynodale}, qui stocke le polygone géospatial
de chacune des 22~régions (type \texttt{MultiPolygon}, SRID~4326)~;
\texttt{District}, rattaché par clé étrangère à une région~; \texttt{Paroisse},
rattachée à un district, avec un champ \texttt{position} de type
\texttt{Point} (SRID~4326) pour les coordonnées GPS, un indicateur
d'activité et les informations de contact~; \texttt{OEuvre}, rattachée à
une paroisse, avec un champ de type d'établissement (école, centre médical,
exploitation agropastorale, immeuble) et un champ de position facultatif.

\noindent\textbf{Ressources humaines.}
L'entité \texttt{Ouvrier} modélise le personnel ecclésiastique~: pasteur,
évangéliste ou diacre permanent, avec leur fonction et leur affectation à
une paroisse. L'entité \texttt{StatistiqueAnnuelle} est associée à une
paroisse et à une année~; elle enregistre les indicateurs démographiques
et financiers annuels.

\noindent\textbf{Gestion de la plateforme.}
L'entité \texttt{CompteUtilisateur} étend le modèle d'utilisateur standard~;
elle porte le rôle, les clés étrangères de périmètre synodal (région,
district ou paroisse), un champ de permissions personnalisées et un
indicateur de première connexion. L'entité \texttt{LogActivite} enregistre
de façon immuable chaque opération d'écriture~: auteur, type d'action,
modèle cible, identifiant de l'objet et horodatage UTC. Les entités visiteur
(\texttt{ProfilVisiteur}, \texttt{ParoisseEnregistree},
\texttt{ItinerairePersonnel}) complètent le schéma pour le module public
authentifié.

\begin{figure}[H]
  \centering
  \fbox{\parbox{0.88\textwidth}{\centering\vspace{3.2cm}
    \textit{[Schéma relationnel-spatial PostGIS
    --- à intégrer après génération draw.io]}%
    \vspace{3.2cm}}}
  \caption{Schéma de la base de données relationnelle-spatiale PostGIS}
  \label{fig:diag-bd}
\end{figure}

%------------------------------------------------------------
\subsection{Diagrammes de Séquence Système}
\label{ssec:seq-systeme}

Les diagrammes de séquence système modélisent les interactions entre les
acteurs externes et le système considéré comme une boîte noire~: ils montrent
les messages échangés et leur ordre chronologique, sans révéler les composants
internes. Deux scénarios représentatifs sont retenus.

\subsubsection{Séquence Système~1~--- Authentification et Accès au Tableau de Bord}

Ce diagramme (Figure~\ref{fig:seq-sys1}) représente le flux nominal de
connexion d'un administrateur~: saisie des identifiants, vérification par le
système, ouverture de session et redirection vers le tableau de bord adapté
au rôle. Il inclut le scénario alternatif de première connexion qui déclenche
obligatoirement le changement de mot de passe avant l'accès aux fonctions
d'administration.

\begin{figure}[H]
  \centering
  \fbox{\parbox{0.88\textwidth}{\centering\vspace{3.5cm}
    \textit{[Diagramme de séquence système 1
    --- Authentification et accès au tableau de bord
    --- à intégrer après génération draw.io]}%
    \vspace{3.5cm}}}
  \caption{Diagramme de séquence système --- Authentification et accès au tableau de bord}
  \label{fig:seq-sys1}
\end{figure}

\subsubsection{Séquence Système~2~--- Consultation Cartographique avec Filtrage}

Ce diagramme (Figure~\ref{fig:seq-sys2}) représente le flux de consultation
de la carte interactive par un utilisateur~: chargement de la carte, affichage
des couches institutionnelles, application d'un filtre hiérarchique par région
et district, mise à jour dynamique des marqueurs et consultation de la fiche
d'une entité sélectionnée. Il illustre la coordination entre l'interface
cartographique et le serveur de données géographiques.

\begin{figure}[H]
  \centering
  \fbox{\parbox{0.88\textwidth}{\centering\vspace{3.5cm}
    \textit{[Diagramme de séquence système 2
    --- Consultation cartographique avec filtrage hiérarchique
    --- à intégrer après génération draw.io]}%
    \vspace{3.5cm}}}
  \caption{Diagramme de séquence système --- Consultation cartographique avec filtrage}
  \label{fig:seq-sys2}
\end{figure}

%------------------------------------------------------------
\subsection{Diagrammes de Séquence Technique}
\label{ssec:seq-technique}

Les diagrammes de séquence technique descendent au niveau des composants
internes du système~: ils montrent les appels entre le frontend, le backend,
la base de données et le serveur cartographique, avec les noms de méthodes
et les flux de données réels.

\subsubsection{Séquence Technique~1~--- Gestion de Session Django}

Ce diagramme (Figure~\ref{fig:seq-tech1}) détaille le traitement interne
d'une requête d'authentification~: la validation des identifiants par le
backend, la création d'une entrée de session dans Redis, la génération du
cookie sécurisé \textit{HttpOnly} et sa transmission au navigateur, ainsi
que la vérification du jeton CSRF sur les requêtes modifiantes suivantes.
Il illustre la chaîne de sécurité complète mise en œuvre entre le frontend
Next.js et le backend Django.

\begin{figure}[H]
  \centering
  \fbox{\parbox{0.88\textwidth}{\centering\vspace{3.5cm}
    \textit{[Diagramme de séquence technique 1
    --- Authentification Django avec gestion de session et CSRF
    --- à intégrer après génération draw.io]}%
    \vspace{3.5cm}}}
  \caption{Diagramme de séquence technique --- Authentification Django et gestion de session}
  \label{fig:seq-tech1}
\end{figure}

\subsubsection{Séquence Technique~2~--- Requête Cartographique Leaflet vers GeoServer}

Ce diagramme (Figure~\ref{fig:seq-tech2}) détaille la chaîne de traitement
d'une requête cartographique émise par Leaflet.js~: construction de la
requête WMS par le client, transmission à GeoServer, interrogation de
PostGIS pour les entités géographiques, application du style SLD, génération
de la tuile image PNG et rendu dans l'interface. Il illustre le rôle de
GeoServer comme couche intermédiaire entre la base de données spatiale et
le client cartographique.

\begin{figure}[H]
  \centering
  \fbox{\parbox{0.88\textwidth}{\centering\vspace{3.5cm}
    \textit{[Diagramme de séquence technique 2
    --- Requête WMS Leaflet → GeoServer → PostGIS
    --- à intégrer après génération draw.io]}%
    \vspace{3.5cm}}}
  \caption{Diagramme de séquence technique --- Requête cartographique WMS de Leaflet vers GeoServer}
  \label{fig:seq-tech2}
\end{figure}

%------------------------------------------------------------
\subsection{Diagramme de Classes Technique}
\label{ssec:classes-tech}

Le diagramme de classes technique (Figure~\ref{fig:diag-classes-tech})
représente le modèle d'implémentation de la plateforme~: il étend le
diagramme de classes métier (Figure~\ref{fig:diag-classes-metier}) avec
les détails d'implémentation propres à Django et à GeoDjango~:
types de données précis (\texttt{PointField}, \texttt{MultiPolygonField},
\texttt{JSONField}), visibilité des attributs, signatures des méthodes
(sérialiseurs DRF, méthodes de filtrage RBAC, helpers de calcul spatial).
Les dépendances entre les classes de modèles, les sérialiseurs et les vues
DRF sont représentées explicitement.

\begin{figure}[H]
  \centering
  \fbox{\parbox{0.88\textwidth}{\centering\vspace{3.2cm}
    \textit{[Diagramme de classes technique
    --- à intégrer après génération draw.io]}%
    \vspace{3.2cm}}}
  \caption{Diagramme de classes technique de l'implémentation Django/GeoDjango}
  \label{fig:diag-classes-tech}
\end{figure}

\medskip
\noindent\textbf{Bilan du chapitre.}
Ce chapitre a formalisé l'ensemble des besoins de la plateforme en
dix-huit exigences fonctionnelles et dix exigences non-fonctionnelles, puis
modélisé le système à travers six diagrammes d'analyse --- contexte,
packages, cas d'utilisation, classes métier et activité --- et cinq
diagrammes de conception --- architecture technique, schéma de données
PostGIS, deux diagrammes de séquence système, deux diagrammes de séquence
technique et diagramme de classes technique. Ces modèles constituent le
contrat de spécification sur lequel s'appuie l'implémentation détaillée au
\textbf{Chapitre~3}.

\clearpage

% ============================================================
%  CHAPITRE 3 — IMPLÉMENTATION ET RÉSULTATS
% ============================================================
\chapter{Implémentation et Résultats}
\label{chap:implementation}

\begin{chapplan}
  \item \textbf{Section~\ref{sec:outils} --- Outils et Technologies}~:
    justification des choix technologiques pour chaque composant de la
    pile logicielle~; tableau récapitulatif des outils, versions et rôles.
  \item \textbf{Section~\ref{sec:deploiement} --- Diagramme de Déploiement}~:
    distribution des cinq services sur les nœuds physiques et virtuels~;
    configuration réseau Docker Compose.
  \item \textbf{Section~\ref{sec:resultats} --- Résultats}~:
    présentation annotée des interfaces fonctionnelles~: carte interactive,
    authentification, tableaux de bord par rôle, import/export et journal
    d'audit~; confrontation aux exigences du Chapitre~2.
  \item \textbf{Section~\ref{sec:discussion} --- Discussion}~:
    vérification des objectifs, analyse critique des forces et des limites
    de la solution réalisée.
\end{chapplan}

\lettrine[lines=2, loversize=0.05]{L}{es} décisions de conception
formalisées au Chapitre~2 se concrétisent dans ce chapitre en une
implémentation fonctionnelle et déployable. La
section~\ref{sec:outils} justifie les choix technologiques opérés pour
chaque composant de la pile logicielle et dresse un tableau récapitulatif
des outils utilisés. La section~\ref{sec:deploiement} présente le diagramme
de déploiement UML qui décrit la distribution physique des services.
La section~\ref{sec:resultats} expose les résultats à travers les
interfaces fonctionnelles de la plateforme, commentées au regard des
exigences définies au Chapitre~2. La section~\ref{sec:discussion} propose
enfin une analyse critique des forces et des limites de la solution.

%------------------------------------------------------------
\section{Outils et Technologies}
\label{sec:outils}
%------------------------------------------------------------

\subsection{Justification des Choix Technologiques}
\label{ssec:justification}

Le choix de chaque technologie résulte d'une évaluation explicite des
alternatives disponibles au regard des contraintes identifiées~: contraintes
institutionnelles (déploiement sur infrastructure limitée, données
sensibles), contraintes fonctionnelles (requêtes spatiales, RBAC
hiérarchique, publication OGC) et contraintes de maintenabilité
(logiciels libres, documentation active, communauté de support).

\subsubsection{Django~5 et Django REST Framework}

Django~5~\cite{django2024} a été retenu comme framework backend principal
pour trois raisons déterminantes. En premier lieu, son extension spatiale
native GeoDjango offre une intégration directe avec PostGIS~: les requêtes
de distance (\texttt{Distance()}), les filtres spatiaux (\texttt{contains},
\texttt{intersects}) et les champs de géométrie (\texttt{PointField},
\texttt{MultiPolygonField}) sont disponibles via l'ORM sans SQL manuel.
En second lieu, le module d'authentification de Django~5 propose une gestion
de sessions robuste par cookies \textit{HttpOnly}~; ce mécanisme est
nativement protégé contre les attaques de type XSS, contrairement aux
jetons JWT transmis dans le stockage local du navigateur, qui présentaient
une vulnérabilité documentée dans la version antérieure de la plateforme.
En troisième lieu, l'écosystème Python disponible ---
\texttt{openpyxl} pour la lecture des fichiers Excel EEC,
\texttt{reportlab} pour la génération PDF, \texttt{django-cors-headers}
pour la politique d'origines croisées --- couvre l'ensemble des besoins
fonctionnels sans dépendances tierces lourdes. Django REST
Framework~\cite{drf2024} complète Django en exposant les ressources
institutionnelles sous forme d'API REST JSON consommée par le frontend.

\subsubsection{PostgreSQL~16 et PostGIS~3.4}

PostgreSQL~16 associé à PostGIS~3.4~\cite{postgis2023} constitue le
standard de facto pour les bases de données spatiales open source~: il est
utilisé dans des projets cartographiques institutionnels à l'échelle nationale
et internationale, dispose d'une indexation spatiale performante
(\textit{R-tree GiST}) et garantit la conformité aux spécifications OGC
Simple Features. La fonction \texttt{ST\_Distance()} native permet
d'identifier la paroisse la plus proche d'une position GPS en une seule
requête SQL, sans traitement applicatif. Aucun SGBD alternatif ne combine
ces fonctionnalités spatiales et une intégration aussi mature avec GeoDjango.

\subsubsection{GeoServer~2.26}

GeoServer~2.26~\cite{geoserver2024} a été retenu pour la publication des
couches cartographiques institutionnelles au regard de sa conformité native
aux protocoles OGC WMS~1.3.0 et WFS~2.0.0, qui garantissent
l'interopérabilité avec tout client cartographique standard. La
configuration des styles visuels via des fichiers SLD permet de définir les
couleurs institutionnelles de l'EEC (vert synodal, fond clair) directement
sur le serveur, sans traitement côté client. L'alternative d'une
publication de fichiers GeoJSON statiques, utilisée dans la version
antérieure, était incompatible avec la mise à jour dynamique des couches
et ne répondait pas aux exigences ENF05 d'interopérabilité OGC.

\subsubsection{Next.js~14}

Next.js~14~\cite{nextjs2024} avec l'App Router a été retenu pour le frontend
pour sa capacité à combiner rendu côté serveur (\textit{SSR}) et rendu côté
client (\textit{CSR}) dans la même application. Le \textit{SSR} est
utilisé pour les pages d'administration~: les données sensibles ne sont
jamais envoyées au navigateur avant vérification de la session, ce qui
réduit la surface d'attaque. Le rendu côté client est utilisé pour la carte
interactive, qui requiert des interactions dynamiques non compatibles avec
le \textit{SSR}. L'App Router de Next.js~14 permet de définir des
\textit{layouts} partagés par rôle d'administration (SUPER/REGION,
DISTRICT, PAROISSE), ce qui simplifie l'implémentation du RBAC côté
frontend.

\subsubsection{Leaflet.js et MarkerCluster}

Leaflet.js~\cite{leaflet2024} a été retenu pour la bibliothèque
cartographique cliente pour sa légèreté (environ 42~Ko minifié), sa licence
libre et l'étendue de son écosystème de plugins. Le plugin
\texttt{Leaflet.MarkerCluster} implémente le regroupement automatique des
565~marqueurs de paroisses en indicateurs numériques~: sans ce mécanisme,
la densité des marqueurs rendrait la carte illisible aux niveaux de zoom
nationaux. La prise en charge native des couches \textit{TileLayer.WMS}
permet de consommer directement les couches GeoServer sans adaptation.

\subsubsection{Redis~7 et Docker Compose}

Redis~7~\cite{redis2024} assure deux fonctions complémentaires~:
le stockage des sessions Django côté serveur (évitant leur persistance
en base de données relationnelle) et la mise en cache des réponses API
fréquentes. Docker Compose~\cite{docker2024} orchestre l'ensemble des
cinq services dans une configuration déclarative reproductible sur toute
machine hôte disposant de Docker~: c'est la réponse directe à l'exigence
ENF06 de portabilité et à la limite de déploiement identifiée dans la
plateforme GeoEEC~v1.0.

%-----
\subsection{Récapitulatif des Outils Utilisés}
\label{ssec:outils-table}

Le Tableau~\ref{tab:outils} dresse la liste complète des outils et
technologies utilisés dans le projet, avec leur version exacte et leur
rôle précis.

\begin{table}[H]
\centering
\caption{Récapitulatif des outils et technologies de la plateforme EEC}
\label{tab:outils}
\setlength{\tabcolsep}{5pt}
\renewcommand{\arraystretch}{1.22}
\resizebox{\linewidth}{!}{%
\begin{tabular}{%
  >{\bfseries}m{3.6cm}
  >{\centering\arraybackslash}m{1.6cm}
  m{9.0cm}}
\toprule
\rowcolor{eecgreen}
\multicolumn{1}{>{\bfseries\color{white}}m{3.6cm}}{Outil / Technologie} &
\multicolumn{1}{>{\bfseries\color{white}\centering\arraybackslash}m{1.6cm}}{Version} &
\multicolumn{1}{>{\bfseries\color{white}}m{9.0cm}}{Rôle dans le projet} \\
\midrule
\rowcolor{eeclightgreen!60}
Python           & 3.12   & Langage principal du backend \\
Django           & 5.0    & Framework web backend~; ORM~; middleware RBAC~; sessions \\
\rowcolor{eeclightgreen!60}
GeoDjango        & 5.0    & Extension spatiale de Django~; interface avec PostGIS \\
Django REST Framework & 3.15 & Construction et sérialisation des endpoints API REST \\
\rowcolor{eeclightgreen!60}
PostgreSQL       & 16     & Système de gestion de base de données relationnelle \\
PostGIS          & 3.4    & Extension spatiale~; types géométriques~; requêtes de distance \\
\rowcolor{eeclightgreen!60}
GeoServer        & 2.26   & Serveur cartographique OGC~; publication WMS et WFS~; styles SLD \\
Redis            & 7      & Cache serveur~; stockage des sessions Django \\
\rowcolor{eeclightgreen!60}
Next.js          & 14     & Framework frontend React~; App Router~; SSR et CSR \\
TypeScript       & 5.x    & Typage statique du code frontend \\
\rowcolor{eeclightgreen!60}
Tailwind CSS     & 3.x    & Framework de mise en forme CSS utilitaire \\
Leaflet.js       & 1.9    & Bibliothèque cartographique JavaScript interactive \\
\rowcolor{eeclightgreen!60}
Leaflet.MarkerCluster & 1.5 & Regroupement automatique des marqueurs sur la carte \\
openpyxl         & 3.1    & Lecture et écriture des fichiers Excel (.xlsx) \\
\rowcolor{eeclightgreen!60}
reportlab        & 4.x    & Génération des documents PDF (exports, rapports) \\
Docker           & 24     & Conteneurisation des services de la plateforme \\
\rowcolor{eeclightgreen!60}
Docker Compose   & 2.x    & Orchestration des cinq services en environnement multi-conteneurs \\
draw.io / diagrams.net & --- & Conception des diagrammes UML du rapport \\
\bottomrule
\end{tabular}}
\end{table}

%------------------------------------------------------------
\section{Diagramme de Déploiement}
\label{sec:deploiement}
%------------------------------------------------------------

Le diagramme de déploiement (Figure~\ref{fig:diag-deploiement}) représente
la distribution physique et virtuelle des composants logiciels de la
plateforme sur les nœuds d'exécution. Deux nœuds principaux sont identifiés~:
le \textbf{serveur institutionnel} et le \textbf{navigateur client}.

Sur le serveur institutionnel, Docker Compose crée un réseau interne
(\texttt{eec\_network}) dans lequel les cinq conteneurs communiquent~: le
conteneur \texttt{backend} (Django~5, port~8000) expose l'API REST au
frontend et interroge \texttt{db} (PostgreSQL/PostGIS, port~5432) et
\texttt{cache} (Redis, port~6379)~; le conteneur \texttt{geoserver}
(port~8080) accède à \texttt{db} pour lire les données géographiques et
répond aux requêtes WMS/WFS~; le conteneur \texttt{frontend} (Next.js~14,
port~3000) sert l'interface web. Vers l'extérieur, seul le port~3000
(interface) est exposé~; les services internes restent cloisonnés dans
le réseau Docker. Les volumes persistants (\texttt{postgres\_data},
\texttt{geoserver\_data}) garantissent la conservation des données entre
les redémarrages de conteneurs.

Sur le navigateur client, Leaflet.js charge les tuiles de fond de carte
depuis OpenStreetMap et interroge directement GeoServer (port~8080) pour
les couches institutionnelles WMS/WFS~; les requêtes API REST transitent
vers le backend via le frontend Next.js.

\begin{figure}[H]
  \centering
  \fbox{\parbox{0.88\textwidth}{\centering\vspace{3.2cm}
    \textit{[Diagramme de déploiement UML
    --- à intégrer après génération draw.io]}%
    \vspace{3.2cm}}}
  \caption{Diagramme de déploiement UML --- Distribution des services Docker Compose}
  \label{fig:diag-deploiement}
\end{figure}

%------------------------------------------------------------
\section{Résultats}
\label{sec:resultats}
%------------------------------------------------------------

Cette section présente les interfaces fonctionnelles de la plateforme
implémentée, commentées au regard des exigences définies au
Chapitre~2. Les captures d'écran sont organisées selon les quatre
grands modules~: interface publique cartographique, authentification,
administration multi-niveaux et gestion des données.

%-----
\subsection{Interface Publique Cartographique}
\label{ssec:res-carte}

L'interface publique est accessible sans authentification depuis la racine
de la plateforme. Elle répond directement aux exigences EF03, EF04, EF05
et EF16 relatives à la cartographie interactive.

La Figure~\ref{fig:res-carte-publique} illustre la carte au niveau de zoom
national~: les 553~paroisses importées sont regroupées en indicateurs de
cluster numériques selon leur densité géographique. Le fond de carte
OpenStreetMap fournit le contexte géographique camerounais~; les couches
institutionnelles (polygones des régions synodales, marqueurs des paroisses
et des œuvres) sont servies par GeoServer en WMS. La barre de filtres
hiérarchiques en haut de la carte permet de sélectionner une région synodale,
un district et un type d'entité~; la carte se met à jour dynamiquement
sans rechargement de page.

\begin{figure}[H]
  \centering
  % ===== IMAGE 1 : Carte publique avec clustering =====
  \fbox{\parbox{0.88\textwidth}{\centering\vspace{3.8cm}
    \textit{[Capture d'écran 1 --- Carte interactive publique avec clustering de marqueurs]}%
    \vspace{3.8cm}}}
  \caption{Interface publique --- Carte interactive avec regroupement automatique des marqueurs}
  \label{fig:res-carte-publique}
\end{figure}

La Figure~\ref{fig:res-popup} montre le panneau de détail qui s'affiche lors
d'un clic sur un marqueur individuel~: nom de la paroisse, district et région
d'appartenance, coordonnées GPS, contacts et statut d'activité. Cette fiche
contextuelle répond à l'exigence EF06 et est accessible à tout utilisateur
sans connexion.

\begin{figure}[H]
  \centering
  % ===== IMAGE 2 : Popup de détail d'une paroisse =====
  \fbox{\parbox{0.88\textwidth}{\centering\vspace{3.8cm}
    \textit{[Capture d'écran 2 --- Popup de détail d'une paroisse sur la carte]}%
    \vspace{3.8cm}}}
  \caption{Interface publique --- Popup de détail d'une paroisse sélectionnée}
  \label{fig:res-popup}
\end{figure}

%-----
\subsection{Interface d'Authentification}
\label{ssec:res-auth}

La Figure~\ref{fig:res-login} présente la page de connexion de l'espace
d'administration (\texttt{/admin/login}). Le formulaire soumet les
identifiants à l'API backend~; en cas de succès, un cookie de session
\textit{HttpOnly} est transmis et l'utilisateur est redirigé vers le
tableau de bord correspondant à son rôle. En cas d'identifiants incorrects,
un message d'erreur est affiché sans révéler si l'identifiant ou le mot de
passe est la source de l'erreur, conformément aux bonnes pratiques de
sécurité recommandées par l'OWASP~\cite{owasp2021}. Pour un compte dont
l'indicateur de première connexion est actif, la redirection s'effectue
automatiquement vers le formulaire de changement de mot de passe.

\begin{figure}[H]
  \centering
  % ===== IMAGE 3 : Page de connexion administrateur =====
  \fbox{\parbox{0.88\textwidth}{\centering\vspace{3.8cm}
    \textit{[Capture d'écran 3 --- Page de connexion de l'espace administrateur]}%
    \vspace{3.8cm}}}
  \caption{Interface d'authentification --- Page de connexion sécurisée}
  \label{fig:res-login}
\end{figure}

%-----
\subsection{Interface d'Administration Multi-Niveaux}
\label{ssec:res-admin}

L'interface d'administration adapte son contenu et ses fonctionnalités au
rôle de l'utilisateur connecté, conformément aux exigences EF02 et EF07 à
EF09. Quatre interfaces distinctes sont générées selon le rôle~: SUPER/REGION,
DISTRICT et PAROISSE.

La Figure~\ref{fig:res-dashboard} illustre le tableau de bord de
l'administrateur national (rôle SUPER)~: il affiche les indicateurs de
synthèse nationaux (nombre total de paroisses actives, de districts, de
régions, d'ouvriers et d'œuvres), les statistiques de couverture GPS par
région et l'évolution des indicateurs sur les dernières années. Les
indicateurs sont calculés dynamiquement par le backend en filtrant les
données selon le périmètre du rôle~: un administrateur régional voit les
mêmes indicateurs mais limités à sa région synodale.

\begin{figure}[H]
  \centering
  % ===== IMAGE 4 : Tableau de bord administrateur national =====
  \fbox{\parbox{0.88\textwidth}{\centering\vspace{3.8cm}
    \textit{[Capture d'écran 4 --- Tableau de bord de l'administrateur national]}%
    \vspace{3.8cm}}}
  \caption{Interface d'administration --- Tableau de bord de l'administrateur national}
  \label{fig:res-dashboard}
\end{figure}

La Figure~\ref{fig:res-paroisses} présente l'interface de gestion des
paroisses~: liste paginée filtrée par région et district, formulaire de
création et de modification avec validation des coordonnées GPS en temps
réel, et action de suppression avec confirmation. Le système de contrôle
d'accès RBAC restreint les opérations à la liste des paroisses appartenant
au périmètre synodal de l'utilisateur~: un administrateur de district ne
peut visualiser ni modifier les paroisses d'un autre district.

\begin{figure}[H]
  \centering
  % ===== IMAGE 5 : Interface de gestion des paroisses =====
  \fbox{\parbox{0.88\textwidth}{\centering\vspace{3.8cm}
    \textit{[Capture d'écran 5 --- Interface de gestion des paroisses (liste et formulaire)]}%
    \vspace{3.8cm}}}
  \caption{Interface d'administration --- Gestion des paroisses avec contrôle RBAC}
  \label{fig:res-paroisses}
\end{figure}

%-----
\subsection{Module d'Import et d'Export}
\label{ssec:res-import}

Le module d'import (Figure~\ref{fig:res-import}) répond à l'exigence EF10~:
il permet de charger les fichiers Excel fournis par les services de l'EEC
directement dans la base de données, après une validation ligne par ligne.
Le rapport d'import affiché après traitement détaille le nombre de lignes
insérées avec succès, le nombre de lignes rejetées et, pour chaque
rejet, la raison précise (coordonnées GPS hors territoire, district
inexistant, doublon détecté). Cette traçabilité de l'import répond
directement au problème de qualité des données identifié lors de l'audit
préliminaire~: 127~paroisses sources sans coordonnées GPS valides ont ainsi
été identifiées et signalées lors de la première importation.

\begin{figure}[H]
  \centering
  % ===== IMAGE 6 : Module d'import Excel et rapport de résultats =====
  \fbox{\parbox{0.88\textwidth}{\centering\vspace{3.8cm}
    \textit{[Capture d'écran 6 --- Module d'import Excel avec rapport détaillé de résultats]}%
    \vspace{3.8cm}}}
  \caption{Module d'import --- Chargement de données Excel et rapport de validation}
  \label{fig:res-import}
\end{figure}

%-----
\subsection{Journal d'Audit}
\label{ssec:res-audit}

Le journal d'audit (Figure~\ref{fig:res-audit}) est accessible
exclusivement à l'administrateur national. Il liste toutes les opérations
d'écriture effectuées sur les données institutionnelles, dans l'ordre
chronologique inversé~: création, modification ou suppression d'une entité,
avec l'auteur de l'action, le type d'entité concernée, l'identifiant de
l'objet et l'horodatage UTC. Des filtres par type d'action, type d'entité
et période permettent de cibler une recherche. En cliquant sur une entrée,
l'administrateur accède au détail de la modification~: valeurs avant et
après la mise à jour. Ce module répond aux exigences EF14 et EF15 relatives
à la traçabilité, et constitue l'un des apports distinctifs par rapport
à la plateforme GeoEEC~v1.0 qui ne disposait d'aucun mécanisme d'audit.

\begin{figure}[H]
  \centering
  % ===== IMAGE 7 : Journal d'audit =====
  \fbox{\parbox{0.88\textwidth}{\centering\vspace{3.8cm}
    \textit{[Capture d'écran 7 --- Journal d'audit des modifications institutionnelles]}%
    \vspace{3.8cm}}}
  \caption{Journal d'audit --- Historique des opérations d'écriture sur les données}
  \label{fig:res-audit}
\end{figure}

%------------------------------------------------------------
\section{Discussion}
\label{sec:discussion}
%------------------------------------------------------------

\subsection{Adéquation aux Objectifs}
\label{ssec:disc-objectifs}

Les six objectifs formulés en Introduction Générale sont confrontés aux
résultats obtenus. L'objectif~1 (modélisation du schéma de données
hiérarchique) est atteint~: le schéma PostGIS implémenté couvre les
22~régions synodales, 137~districts, 553~paroisses importées, 311~œuvres
et 685~ouvriers, avec toutes les contraintes d'intégrité référentielle et
les champs de géolocalisation. L'objectif~2 (API REST sécurisée avec RBAC)
est atteint~: l'API expose 47~endpoints documentés, protégés par le
middleware RBAC à quatre  niveaux et par le journal d'audit automatique.
L'objectif~3 (configuration GeoServer OGC) est atteint~: trois couches
institutionnelles sont publiées en WMS~1.3.0 et WFS~2.0.0 avec styles SLD.
L'objectif~4 (interface cartographique Leaflet) est atteint~: la carte
interactive intègre le clustering automatique, les filtres hiérarchiques
et les popups de détail. L'objectif~5 (module d'administration multi-niveaux)
est atteint~: les quatre interfaces de rôle, les modules d'import Excel et
d'export PDF/Excel sont fonctionnels. L'objectif~6 (conteneurisation Docker)
est atteint~: la plateforme est entièrement déployable avec une commande
unique sur tout serveur disposant de Docker.

%-----
\subsection{Forces de la Solution}
\label{ssec:disc-forces}

\textbf{Sécurité améliorée par rapport à la version antérieure.}
Le remplacement de l'authentification JWT par des sessions Django avec
cookies \textit{HttpOnly} et \textit{SameSite=Strict} élimine le vecteur
d'attaque XSS qui permettait la capture des jetons transmis dans le stockage
local du navigateur. La protection CSRF systématique sur toutes les
requêtes modifiantes renforce ce dispositif, conformément aux recommandations
de l'OWASP~\cite{owasp2021}.

\textbf{Interopérabilité cartographique.}
La publication des couches géographiques via GeoServer en WMS~1.3.0 et
WFS~2.0.0 rend les données de l'EEC consommables par tout client
cartographique compatible OGC~: QGIS, OpenLayers, ou un futur portail
institutionnel camerounais. Aucune solution concurrente analysée au
Chapitre~1 ne proposait cette combinaison.

\textbf{RBAC aligné sur la hiérarchie réelle de l'EEC.}
Le modèle de contrôle d'accès à quatre niveaux --- national, régional,
district et paroissial --- correspond exactement à la structure synodale
de l'institution~: un administrateur de district peut gérer les paroisses
de son périmètre sans accès aux données des districts voisins. Ce niveau
de granularité était absent de la version antérieure.

\textbf{Traçabilité complète et automatique.}
Tout enregistrement de modification, de création ou de suppression
déclenche automatiquement l'écriture d'une entrée dans le journal d'audit
via les signaux Django, sans intervention du développeur dans chaque vue.
Cette automatisation garantit l'exhaustivité de l'historique, impossible
à assurer manuellement.

\textbf{Déploiement reproductible.}
La configuration Docker Compose en cinq services garantit que l'environnement
d'exécution est identique en développement et en production~: plus aucune
divergence de versions de bibliothèques ou de configuration système ne peut
causer des défaillances à l'installation.

Ces cinq forces confirment la pertinence des décisions architecturales
formalisées au Chapitre~2. Elles positionnent la plateforme de façon
distincte par rapport aux solutions analysées dans l'état de l'art, en
particulier sur les axes de sécurité, d'interopérabilité OGC et de
RBAC institutionnel. La section suivante expose les limites structurelles
qui tempèrent la portée de ces résultats.

%------------------------------------------------------------
\section{Limites et Contraintes}
\label{sec:limites}
%------------------------------------------------------------

Tout travail d'ingénierie comporte des contraintes qui en délimitent la
portée. Une évaluation honnête de ces limites est indispensable pour
apprécier correctement la valeur de la solution et orienter les
développements futurs.

\textbf{Couverture GPS incomplète.}
L'audit des données sources a révélé que 127~paroisses sur 565, soit
22~\% de l'effectif total, ne disposent d'aucune coordonnée GPS valide
dans les fichiers Excel fournis par l'EEC. Ces paroisses sont importées
en base de données mais ne s'affichent pas sur la carte interactive.
La plateforme ne peut donc pas prétendre à une couverture cartographique
complète du réseau ecclésiastique~; son utilité pour les décisions de
planification territoriale nationale reste limitée tant que ce taux
n'est pas réduit par une campagne de collecte terrain.

\textbf{Distance à vol d'oiseau.}
La fonction de calcul d'itinéraire repose sur la formule de Haversine,
qui mesure la distance géodésique en ligne droite entre deux points.
Cette mesure sous-estime systématiquement la durée réelle de trajet
sur un réseau routier camerounais caractérisé par des voies sinueuses
et un relief accidenté~; elle ne tient pas compte des barrières
physiques (fleuves, massifs montagneux, absence de pont).
La durée estimée à 60~km/h de moyenne est donc approximative et ne
peut servir de base à des décisions opérationnelles précises.

\textbf{Absence d'application mobile.}
La plateforme est conçue pour des écrans desktop et tablette
(résolution minimale 1\,024~$\times$~768~px, ENF09). Elle n'est pas
adaptée aux smartphones, qui constituent pourtant le principal outil
numérique des administrateurs paroissiaux en zone rurale camerounaise.
En l'absence d'application mobile dédiée, la collecte terrain des
coordonnées GPS des 127~paroisses non géolocalisées reste sans
outillage numérique adapté.

\textbf{Authentification à facteur unique.}
L'accès à l'espace d'administration repose sur un couple
identifiant/mot de passe sans mécanisme de double authentification
(TOTP ou code temporaire par messagerie). Pour des comptes de niveau
national disposant d'un accès complet à l'ensemble des données de
l'EEC, cette limite expose la plateforme aux risques de compromission
par capture ou réutilisation de mot de passe, en l'absence de second
facteur de vérification.

\textbf{Déploiement mono-instance.}
L'architecture Docker Compose actuelle déploie chaque service en
instance unique sur un seul serveur physique. Elle ne prévoit pas de
mécanisme de réplication ou de bascule automatique en cas de défaillance
matérielle~; la disponibilité de la plateforme est donc directement
dépendante de la fiabilité du serveur institutionnel hôte.

\medskip
\noindent\textbf{Bilan du chapitre.}
Ce chapitre a présenté et justifié les choix technologiques de la
plateforme à travers quatre tableaux comparatifs, exposé les résultats
à travers sept interfaces fonctionnelles commentées, analysé les cinq
forces distinctives de la solution dans la section Discussion, puis
identifié cinq limites structurelles dans cette section. La
\textbf{Conclusion Générale} synthétise l'ensemble de la démarche
et formule les perspectives d'amélioration directement dérivées
de ces limites.

\clearpage

% ============================================================
%  CONCLUSION GÉNÉRALE
% ============================================================
\chapter*{Conclusion Générale}
\addcontentsline{toc}{chapter}{Conclusion Générale}
\markboth{Conclusion Générale}{Conclusion Générale}

\noindent
L'Église Évangélique du Cameroun gérait jusqu'alors ses 565~paroisses,
311~œuvres sociales et 685~ouvriers ecclésiastiques répartis sur
22~régions synodales et 137~districts sans aucun outil numérique
centralisé~: pas de visualisation géographique, pas de contrôle d'accès
formalisé selon les périmètres synodaux, pas de journal d'audit des
modifications, pas de publication cartographique interopérable. Ce travail
a abordé ce problème par la conception et le développement d'une
plateforme web cartographique institutionnelle, en s'appuyant sur
l'hypothèse qu'une architecture à services découplés, entièrement
conteneurisée, combinant un serveur cartographique OGC et un contrôle
d'accès hiérarchique aligné sur la structure synodale, permettrait d'y
répondre de façon complète et déployable dans un contexte institutionnel
à ressources limitées.

Cette hypothèse est validée par les résultats obtenus~: la plateforme
implémentée déploie en une seule commande cinq services
--- Django~5~+~GeoDjango, PostgreSQL~+~PostGIS~3.4, GeoServer~2.26,
Next.js~14 et Redis~7~--- et offre une carte interactive avec clustering
de 553~paroisses, des filtres hiérarchiques par région et district, un
tableau de bord adaptatif par rôle, un module d'import/export des données
institutionnelles, et un journal d'audit automatique et exhaustif de toutes
les opérations d'écriture. L'ensemble des six objectifs formulés en
Introduction Générale a été atteint. Par rapport à la plateforme
GeoEEC~v1.0, la solution apporte une valeur ajoutée précise et documentée
sur cinq axes~: sécurité de l'authentification par sessions \textit{HttpOnly},
RBAC aligné sur les quatre niveaux de la hiérarchie synodale réelle,
publication cartographique conforme aux standards OGC WMS et WFS, traçabilité
automatique par journal d'audit, et reproductibilité de l'environnement
d'exécution par conteneurisation complète.

\medskip
\noindent\textbf{Limites.}
Cinq limites structurelles, analysées en détail dans la
section~\ref{sec:limites} du Chapitre~3, tempèrent la portée de ces
résultats. La couverture cartographique reste incomplète~: 127~paroisses
(22~\%) ne disposent d'aucune coordonnée GPS dans les fichiers sources de
l'EEC et ne s'affichent pas sur la carte, ce qui restreint directement
l'utilité de la plateforme pour les décisions de planification territoriale.
Le calcul d'itinéraire repose sur la distance à vol d'oiseau et non sur
le réseau routier réel, produisant des durées de trajet sous-estimées dans
les zones à relief accidenté. L'absence d'interface mobile prive les
administrateurs paroissiaux d'un outil de saisie GPS sur le terrain.
L'authentification à facteur unique expose les comptes nationaux à des
risques de compromission par mot de passe seul. Enfin, le déploiement
mono-instance ne prévoit pas de mécanisme de haute disponibilité en cas
de défaillance du serveur hôte.

\medskip
\noindent\textbf{Perspectives.}
Ces limites définissent directement les axes de développement prioritaires.
La première perspective est le développement d'une application mobile
légère (React Native ou Progressive Web App) permettant aux administrateurs
paroissiaux de saisir ou de corriger les coordonnées GPS directement sur le
terrain~: cette seule fonctionnalité permettrait de résorber le déficit de
géolocalisation des 127~paroisses non couvertes et de faire passer le taux
de couverture cartographique de 78~\% à 100~\%. La deuxième perspective est
l'intégration d'un moteur de calcul d'itinéraires routiers~--- via l'API
OSRM (\textit{Open Source Routing Machine}) ou une instance locale de
GraphHopper~--- pour remplacer la distance à vol d'oiseau par des temps de
trajet réalistes sur les routes camerounaises. La troisième perspective est
l'activation de la double authentification par TOTP (\textit{Time-based
One-Time Password}) pour les comptes d'administrateurs nationaux et
régionaux, en s'appuyant sur la bibliothèque \texttt{django-otp}
déjà compatible avec l'infrastructure Django. Enfin, la réintégration
du module de questions-réponses en langage naturel développé dans
GeoEEC~v1.0 --- enrichi des nouvelles entités et des statistiques annuelles
--- ouvrirait la plateforme à des usages analytiques avancés permettant
aux responsables synodaux d'interroger les données institutionnelles sans
compétences techniques.

\noindent\rule{\linewidth}{0.8pt}
 dans complet.tex,
%            juste après le contenu de la Section 2.2 existante.
% ============================================================

% ============================================================
%  SUITE — SECTION 2.2 CONCEPTION DU SYSTÈME
% ============================================================

\subsection{Architecture Technique}
\label{ssec:archi}

L'architecture retenue pour la plateforme EEC~Géolocalisation est une
\textbf{architecture à services découplés}, dans laquelle chaque composant
fonctionnel est isolé dans un conteneur indépendant et communique avec les
autres via des interfaces standardisées. Ce choix s'oppose à une architecture
monolithique dans laquelle toutes les fonctions seraient regroupées dans un
seul processus~: la séparation des responsabilités facilite la maintenance,
permet la mise à jour indépendante de chaque couche et renforce la sécurité en
cloisonnant les surfaces d'attaque.

La plateforme est décomposée en cinq services orchestrés par Docker~Compose,
représentés dans la Figure~\ref{fig:diag-archi}~:

\begin{itemize}[itemsep=3pt, topsep=5pt, leftmargin=1.5cm]
  \item \textbf{Backend} --- serveur Django~5 avec GeoDjango et
    Django REST Framework~: expose l'intégralité de la logique métier
    sous forme d'API REST, applique le contrôle d'accès basé sur les rôles
    et gère les sessions d'authentification par cookies sécurisés
    (\textit{HttpOnly}, \textit{SameSite=Strict})~;
  \item \textbf{Base de données} --- PostgreSQL~16 avec l'extension
    PostGIS~3.4~: assure la persistance des données institutionnelles et
    des géométries spatiales (points, polygones)~; les requêtes de
    proximité exploitent les fonctions de calcul de distance natives~;
  \item \textbf{Serveur cartographique} --- GeoServer~2.26~: publie les
    couches géographiques de l'EEC selon les standards OGC WMS~1.3.0 et
    WFS~2.0.0~; les styles visuels sont définis via des fichiers SLD aux
    couleurs institutionnelles~;
  \item \textbf{Cache} --- Redis~7~: stocke les sessions Django côté
    serveur et met en cache les réponses fréquentes de l'API pour réduire
    la charge sur la base de données~;
  \item \textbf{Frontend} --- Next.js~14 avec l'App Router~: sert
    l'interface publique (carte Leaflet, clustering, filtres) et
    l'interface d'administration multi-niveaux~; les pages critiques sont
    rendues côté serveur (\textit{SSR}) pour garantir la sécurité des
    données sensibles.
\end{itemize}

Les échanges entre le frontend et le backend transitent exclusivement par
l'API REST en JSON. Le frontend Leaflet interroge GeoServer directement
pour les couches cartographiques, sans passer par le backend Django~: cette
séparation permet d'exploiter les protocoles OGC standard sans surcharger
le serveur applicatif. L'ensemble des services partage un réseau Docker
interne isolé~; seuls les ports nécessaires sont exposés vers l'hôte.

\begin{figure}[H]
  \centering
  \fbox{\parbox{0.88\textwidth}{\centering\vspace{3.2cm}
    \textit{[Diagramme d'architecture technique à services découplés
    --- à intégrer après export draw.io]}%
    \vspace{3.2cm}}}
  \caption{Architecture technique à services découplés de la plateforme EEC}
  \label{fig:diag-archi}
\end{figure}

%------------------------------------------------------------
\subsection{Modèle de Données}
\label{ssec:modele-bd}

Le schéma de données de la plateforme est un modèle relationnel-spatial
conforme aux normes PostGIS. Il est organisé en trois groupes d'entités
directement dérivés du domaine institutionnel de l'EEC et des exigences
fonctionnelles identifiées à la section~\ref{ssec:exigences}.

\noindent\textbf{Entités géographiques hiérarchiques.}
La structure territoriale de l'EEC est modélisée par quatre tables liées
en cascade~: \texttt{RegionSynodale}, qui stocke le polygone géospatial
de chacune des 22~régions (type \texttt{MultiPolygon}, SRID~4326)~;
\texttt{District}, rattaché par clé étrangère à une région~; \texttt{Paroisse},
rattachée à un district, avec un champ \texttt{position} de type
\texttt{Point} (SRID~4326) pour les coordonnées GPS, un indicateur
d'activité et les informations de contact~; \texttt{OEuvre}, rattachée à
une paroisse, avec un champ de type d'établissement (école, centre médical,
exploitation agropastorale, immeuble) et un champ de position facultatif.

\noindent\textbf{Ressources humaines.}
L'entité \texttt{Ouvrier} modélise le personnel ecclésiastique~: pasteur,
évangéliste ou diacre permanent, avec leur fonction et leur affectation à
une paroisse. L'entité \texttt{StatistiqueAnnuelle} est associée à une
paroisse et à une année~; elle enregistre les indicateurs démographiques
et financiers annuels.

\noindent\textbf{Gestion de la plateforme.}
L'entité \texttt{CompteUtilisateur} étend le modèle d'utilisateur standard~;
elle porte le rôle, les clés étrangères de périmètre synodal (région,
district ou paroisse), un champ de permissions personnalisées et un
indicateur de première connexion. L'entité \texttt{LogActivite} enregistre
de façon immuable chaque opération d'écriture~: auteur, type d'action,
modèle cible, identifiant de l'objet et horodatage UTC. Les entités visiteur
(\texttt{ProfilVisiteur}, \texttt{ParoisseEnregistree},
\texttt{ItinerairePersonnel}) complètent le schéma pour le module public
authentifié.

\begin{figure}[H]
  \centering
  \fbox{\parbox{0.88\textwidth}{\centering\vspace{3.2cm}
    \textit{[Schéma relationnel-spatial PostGIS
    --- à intégrer après génération draw.io]}%
    \vspace{3.2cm}}}
  \caption{Schéma de la base de données relationnelle-spatiale PostGIS}
  \label{fig:diag-bd}
\end{figure}

%------------------------------------------------------------
\subsection{Diagrammes de Séquence Système}
\label{ssec:seq-systeme}

Les diagrammes de séquence système modélisent les interactions entre les
acteurs externes et le système considéré comme une boîte noire~: ils montrent
les messages échangés et leur ordre chronologique, sans révéler les composants
internes. Deux scénarios représentatifs sont retenus.

\subsubsection{Séquence Système~1~--- Authentification et Accès au Tableau de Bord}

Ce diagramme (Figure~\ref{fig:seq-sys1}) représente le flux nominal de
connexion d'un administrateur~: saisie des identifiants, vérification par le
système, ouverture de session et redirection vers le tableau de bord adapté
au rôle. Il inclut le scénario alternatif de première connexion qui déclenche
obligatoirement le changement de mot de passe avant l'accès aux fonctions
d'administration.

\begin{figure}[H]
  \centering
  \fbox{\parbox{0.88\textwidth}{\centering\vspace{3.5cm}
    \textit{[Diagramme de séquence système 1
    --- Authentification et accès au tableau de bord
    --- à intégrer après génération draw.io]}%
    \vspace{3.5cm}}}
  \caption{Diagramme de séquence système --- Authentification et accès au tableau de bord}
  \label{fig:seq-sys1}
\end{figure}

\subsubsection{Séquence Système~2~--- Consultation Cartographique avec Filtrage}

Ce diagramme (Figure~\ref{fig:seq-sys2}) représente le flux de consultation
de la carte interactive par un utilisateur~: chargement de la carte, affichage
des couches institutionnelles, application d'un filtre hiérarchique par région
et district, mise à jour dynamique des marqueurs et consultation de la fiche
d'une entité sélectionnée. Il illustre la coordination entre l'interface
cartographique et le serveur de données géographiques.

\begin{figure}[H]
  \centering
  \fbox{\parbox{0.88\textwidth}{\centering\vspace{3.5cm}
    \textit{[Diagramme de séquence système 2
    --- Consultation cartographique avec filtrage hiérarchique
    --- à intégrer après génération draw.io]}%
    \vspace{3.5cm}}}
  \caption{Diagramme de séquence système --- Consultation cartographique avec filtrage}
  \label{fig:seq-sys2}
\end{figure}

%------------------------------------------------------------
\subsection{Diagrammes de Séquence Technique}
\label{ssec:seq-technique}

Les diagrammes de séquence technique descendent au niveau des composants
internes du système~: ils montrent les appels entre le frontend, le backend,
la base de données et le serveur cartographique, avec les noms de méthodes
et les flux de données réels.

\subsubsection{Séquence Technique~1~--- Gestion de Session Django}

Ce diagramme (Figure~\ref{fig:seq-tech1}) détaille le traitement interne
d'une requête d'authentification~: la validation des identifiants par le
backend, la création d'une entrée de session dans Redis, la génération du
cookie sécurisé \textit{HttpOnly} et sa transmission au navigateur, ainsi
que la vérification du jeton CSRF sur les requêtes modifiantes suivantes.
Il illustre la chaîne de sécurité complète mise en œuvre entre le frontend
Next.js et le backend Django.

\begin{figure}[H]
  \centering
  \fbox{\parbox{0.88\textwidth}{\centering\vspace{3.5cm}
    \textit{[Diagramme de séquence technique 1
    --- Authentification Django avec gestion de session et CSRF
    --- à intégrer après génération draw.io]}%
    \vspace{3.5cm}}}
  \caption{Diagramme de séquence technique --- Authentification Django et gestion de session}
  \label{fig:seq-tech1}
\end{figure}

\subsubsection{Séquence Technique~2~--- Requête Cartographique Leaflet vers GeoServer}

Ce diagramme (Figure~\ref{fig:seq-tech2}) détaille la chaîne de traitement
d'une requête cartographique émise par Leaflet.js~: construction de la
requête WMS par le client, transmission à GeoServer, interrogation de
PostGIS pour les entités géographiques, application du style SLD, génération
de la tuile image PNG et rendu dans l'interface. Il illustre le rôle de
GeoServer comme couche intermédiaire entre la base de données spatiale et
le client cartographique.

\begin{figure}[H]
  \centering
  \fbox{\parbox{0.88\textwidth}{\centering\vspace{3.5cm}
    \textit{[Diagramme de séquence technique 2
    --- Requête WMS Leaflet → GeoServer → PostGIS
    --- à intégrer après génération draw.io]}%
    \vspace{3.5cm}}}
  \caption{Diagramme de séquence technique --- Requête cartographique WMS de Leaflet vers GeoServer}
  \label{fig:seq-tech2}
\end{figure}

%------------------------------------------------------------
\subsection{Diagramme de Classes Technique}
\label{ssec:classes-tech}

Le diagramme de classes technique (Figure~\ref{fig:diag-classes-tech})
représente le modèle d'implémentation de la plateforme~: il étend le
diagramme de classes métier (Figure~\ref{fig:diag-classes-metier}) avec
les détails d'implémentation propres à Django et à GeoDjango~:
types de données précis (\texttt{PointField}, \texttt{MultiPolygonField},
\texttt{JSONField}), visibilité des attributs, signatures des méthodes
(sérialiseurs DRF, méthodes de filtrage RBAC, helpers de calcul spatial).
Les dépendances entre les classes de modèles, les sérialiseurs et les vues
DRF sont représentées explicitement.

\begin{figure}[H]
  \centering
  \fbox{\parbox{0.88\textwidth}{\centering\vspace{3.2cm}
    \textit{[Diagramme de classes technique
    --- à intégrer après génération draw.io]}%
    \vspace{3.2cm}}}
  \caption{Diagramme de classes technique de l'implémentation Django/GeoDjango}
  \label{fig:diag-classes-tech}
\end{figure}

\medskip
\noindent\textbf{Bilan du chapitre.}
Ce chapitre a formalisé l'ensemble des besoins de la plateforme en
dix-huit exigences fonctionnelles et dix exigences non-fonctionnelles, puis
modélisé le système à travers six diagrammes d'analyse --- contexte,
packages, cas d'utilisation, classes métier et activité --- et cinq
diagrammes de conception --- architecture technique, schéma de données
PostGIS, deux diagrammes de séquence système, deux diagrammes de séquence
technique et diagramme de classes technique. Ces modèles constituent le
contrat de spécification sur lequel s'appuie l'implémentation détaillée au
\textbf{Chapitre~3}.

\clearpage

% ============================================================
%  CHAPITRE 3 — IMPLÉMENTATION ET RÉSULTATS
% ============================================================
\chapter{Implémentation et Résultats}
\label{chap:implementation}

\begin{chapplan}
  \item \textbf{Section~\ref{sec:outils} --- Outils et Technologies}~:
    justification des choix technologiques pour chaque composant de la
    pile logicielle~; tableau récapitulatif des outils, versions et rôles.
  \item \textbf{Section~\ref{sec:deploiement} --- Diagramme de Déploiement}~:
    distribution des cinq services sur les nœuds physiques et virtuels~;
    configuration réseau Docker Compose.
  \item \textbf{Section~\ref{sec:resultats} --- Résultats}~:
    présentation annotée des interfaces fonctionnelles~: carte interactive,
    authentification, tableaux de bord par rôle, import/export et journal
    d'audit~; confrontation aux exigences du Chapitre~2.
  \item \textbf{Section~\ref{sec:discussion} --- Discussion}~:
    vérification des objectifs, analyse critique des forces et des limites
    de la solution réalisée.
\end{chapplan}

\lettrine[lines=2, loversize=0.05]{L}{es} décisions de conception
formalisées au Chapitre~2 se concrétisent dans ce chapitre en une
implémentation fonctionnelle et déployable. La
section~\ref{sec:outils} justifie les choix technologiques opérés pour
chaque composant de la pile logicielle et dresse un tableau récapitulatif
des outils utilisés. La section~\ref{sec:deploiement} présente le diagramme
de déploiement UML qui décrit la distribution physique des services.
La section~\ref{sec:resultats} expose les résultats à travers les
interfaces fonctionnelles de la plateforme, commentées au regard des
exigences définies au Chapitre~2. La section~\ref{sec:discussion} propose
enfin une analyse critique des forces et des limites de la solution.

%------------------------------------------------------------
\section{Outils et Technologies}
\label{sec:outils}
%------------------------------------------------------------

\subsection{Justification des Choix Technologiques}
\label{ssec:justification}

Le choix de chaque technologie résulte d'une évaluation explicite des
alternatives disponibles au regard des contraintes identifiées~: contraintes
institutionnelles (déploiement sur infrastructure limitée, données
sensibles), contraintes fonctionnelles (requêtes spatiales, RBAC
hiérarchique, publication OGC) et contraintes de maintenabilité
(logiciels libres, documentation active, communauté de support).

\subsubsection{Django~5 et Django REST Framework}

Django~5~\cite{django2024} a été retenu comme framework backend principal
pour trois raisons déterminantes. En premier lieu, son extension spatiale
native GeoDjango offre une intégration directe avec PostGIS~: les requêtes
de distance (\texttt{Distance()}), les filtres spatiaux (\texttt{contains},
\texttt{intersects}) et les champs de géométrie (\texttt{PointField},
\texttt{MultiPolygonField}) sont disponibles via l'ORM sans SQL manuel.
En second lieu, le module d'authentification de Django~5 propose une gestion
de sessions robuste par cookies \textit{HttpOnly}~; ce mécanisme est
nativement protégé contre les attaques de type XSS, contrairement aux
jetons JWT transmis dans le stockage local du navigateur, qui présentaient
une vulnérabilité documentée dans la version antérieure de la plateforme.
En troisième lieu, l'écosystème Python disponible ---
\texttt{openpyxl} pour la lecture des fichiers Excel EEC,
\texttt{reportlab} pour la génération PDF, \texttt{django-cors-headers}
pour la politique d'origines croisées --- couvre l'ensemble des besoins
fonctionnels sans dépendances tierces lourdes. Django REST
Framework~\cite{drf2024} complète Django en exposant les ressources
institutionnelles sous forme d'API REST JSON consommée par le frontend.

\subsubsection{PostgreSQL~16 et PostGIS~3.4}

PostgreSQL~16 associé à PostGIS~3.4~\cite{postgis2023} constitue le
standard de facto pour les bases de données spatiales open source~: il est
utilisé dans des projets cartographiques institutionnels à l'échelle nationale
et internationale, dispose d'une indexation spatiale performante
(\textit{R-tree GiST}) et garantit la conformité aux spécifications OGC
Simple Features. La fonction \texttt{ST\_Distance()} native permet
d'identifier la paroisse la plus proche d'une position GPS en une seule
requête SQL, sans traitement applicatif. Aucun SGBD alternatif ne combine
ces fonctionnalités spatiales et une intégration aussi mature avec GeoDjango.

\subsubsection{GeoServer~2.26}

GeoServer~2.26~\cite{geoserver2024} a été retenu pour la publication des
couches cartographiques institutionnelles au regard de sa conformité native
aux protocoles OGC WMS~1.3.0 et WFS~2.0.0, qui garantissent
l'interopérabilité avec tout client cartographique standard. La
configuration des styles visuels via des fichiers SLD permet de définir les
couleurs institutionnelles de l'EEC (vert synodal, fond clair) directement
sur le serveur, sans traitement côté client. L'alternative d'une
publication de fichiers GeoJSON statiques, utilisée dans la version
antérieure, était incompatible avec la mise à jour dynamique des couches
et ne répondait pas aux exigences ENF05 d'interopérabilité OGC.

\subsubsection{Next.js~14}

Next.js~14~\cite{nextjs2024} avec l'App Router a été retenu pour le frontend
pour sa capacité à combiner rendu côté serveur (\textit{SSR}) et rendu côté
client (\textit{CSR}) dans la même application. Le \textit{SSR} est
utilisé pour les pages d'administration~: les données sensibles ne sont
jamais envoyées au navigateur avant vérification de la session, ce qui
réduit la surface d'attaque. Le rendu côté client est utilisé pour la carte
interactive, qui requiert des interactions dynamiques non compatibles avec
le \textit{SSR}. L'App Router de Next.js~14 permet de définir des
\textit{layouts} partagés par rôle d'administration (SUPER/REGION,
DISTRICT, PAROISSE), ce qui simplifie l'implémentation du RBAC côté
frontend.

\subsubsection{Leaflet.js et MarkerCluster}

Leaflet.js~\cite{leaflet2024} a été retenu pour la bibliothèque
cartographique cliente pour sa légèreté (environ 42~Ko minifié), sa licence
libre et l'étendue de son écosystème de plugins. Le plugin
\texttt{Leaflet.MarkerCluster} implémente le regroupement automatique des
565~marqueurs de paroisses en indicateurs numériques~: sans ce mécanisme,
la densité des marqueurs rendrait la carte illisible aux niveaux de zoom
nationaux. La prise en charge native des couches \textit{TileLayer.WMS}
permet de consommer directement les couches GeoServer sans adaptation.

\subsubsection{Redis~7 et Docker Compose}

Redis~7~\cite{redis2024} assure deux fonctions complémentaires~:
le stockage des sessions Django côté serveur (évitant leur persistance
en base de données relationnelle) et la mise en cache des réponses API
fréquentes. Docker Compose~\cite{docker2024} orchestre l'ensemble des
cinq services dans une configuration déclarative reproductible sur toute
machine hôte disposant de Docker~: c'est la réponse directe à l'exigence
ENF06 de portabilité et à la limite de déploiement identifiée dans la
plateforme GeoEEC~v1.0.

%-----
\subsection{Récapitulatif des Outils Utilisés}
\label{ssec:outils-table}

Le Tableau~\ref{tab:outils} dresse la liste complète des outils et
technologies utilisés dans le projet, avec leur version exacte et leur
rôle précis.

\begin{table}[H]
\centering
\caption{Récapitulatif des outils et technologies de la plateforme EEC}
\label{tab:outils}
\setlength{\tabcolsep}{5pt}
\renewcommand{\arraystretch}{1.22}
\resizebox{\linewidth}{!}{%
\begin{tabular}{%
  >{\bfseries}m{3.6cm}
  >{\centering\arraybackslash}m{1.6cm}
  m{9.0cm}}
\toprule
\rowcolor{eecgreen}
\multicolumn{1}{>{\bfseries\color{white}}m{3.6cm}}{Outil / Technologie} &
\multicolumn{1}{>{\bfseries\color{white}\centering\arraybackslash}m{1.6cm}}{Version} &
\multicolumn{1}{>{\bfseries\color{white}}m{9.0cm}}{Rôle dans le projet} \\
\midrule
\rowcolor{eeclightgreen!60}
Python           & 3.12   & Langage principal du backend \\
Django           & 5.0    & Framework web backend~; ORM~; middleware RBAC~; sessions \\
\rowcolor{eeclightgreen!60}
GeoDjango        & 5.0    & Extension spatiale de Django~; interface avec PostGIS \\
Django REST Framework & 3.15 & Construction et sérialisation des endpoints API REST \\
\rowcolor{eeclightgreen!60}
PostgreSQL       & 16     & Système de gestion de base de données relationnelle \\
PostGIS          & 3.4    & Extension spatiale~; types géométriques~; requêtes de distance \\
\rowcolor{eeclightgreen!60}
GeoServer        & 2.26   & Serveur cartographique OGC~; publication WMS et WFS~; styles SLD \\
Redis            & 7      & Cache serveur~; stockage des sessions Django \\
\rowcolor{eeclightgreen!60}
Next.js          & 14     & Framework frontend React~; App Router~; SSR et CSR \\
TypeScript       & 5.x    & Typage statique du code frontend \\
\rowcolor{eeclightgreen!60}
Tailwind CSS     & 3.x    & Framework de mise en forme CSS utilitaire \\
Leaflet.js       & 1.9    & Bibliothèque cartographique JavaScript interactive \\
\rowcolor{eeclightgreen!60}
Leaflet.MarkerCluster & 1.5 & Regroupement automatique des marqueurs sur la carte \\
openpyxl         & 3.1    & Lecture et écriture des fichiers Excel (.xlsx) \\
\rowcolor{eeclightgreen!60}
reportlab        & 4.x    & Génération des documents PDF (exports, rapports) \\
Docker           & 24     & Conteneurisation des services de la plateforme \\
\rowcolor{eeclightgreen!60}
Docker Compose   & 2.x    & Orchestration des cinq services en environnement multi-conteneurs \\
draw.io / diagrams.net & --- & Conception des diagrammes UML du rapport \\
\bottomrule
\end{tabular}}
\end{table}

%------------------------------------------------------------
\section{Diagramme de Déploiement}
\label{sec:deploiement}
%------------------------------------------------------------

Le diagramme de déploiement (Figure~\ref{fig:diag-deploiement}) représente
la distribution physique et virtuelle des composants logiciels de la
plateforme sur les nœuds d'exécution. Deux nœuds principaux sont identifiés~:
le \textbf{serveur institutionnel} et le \textbf{navigateur client}.

Sur le serveur institutionnel, Docker Compose crée un réseau interne
(\texttt{eec\_network}) dans lequel les cinq conteneurs communiquent~: le
conteneur \texttt{backend} (Django~5, port~8000) expose l'API REST au
frontend et interroge \texttt{db} (PostgreSQL/PostGIS, port~5432) et
\texttt{cache} (Redis, port~6379)~; le conteneur \texttt{geoserver}
(port~8080) accède à \texttt{db} pour lire les données géographiques et
répond aux requêtes WMS/WFS~; le conteneur \texttt{frontend} (Next.js~14,
port~3000) sert l'interface web. Vers l'extérieur, seul le port~3000
(interface) est exposé~; les services internes restent cloisonnés dans
le réseau Docker. Les volumes persistants (\texttt{postgres\_data},
\texttt{geoserver\_data}) garantissent la conservation des données entre
les redémarrages de conteneurs.

Sur le navigateur client, Leaflet.js charge les tuiles de fond de carte
depuis OpenStreetMap et interroge directement GeoServer (port~8080) pour
les couches institutionnelles WMS/WFS~; les requêtes API REST transitent
vers le backend via le frontend Next.js.

\begin{figure}[H]
  \centering
  \fbox{\parbox{0.88\textwidth}{\centering\vspace{3.2cm}
    \textit{[Diagramme de déploiement UML
    --- à intégrer après génération draw.io]}%
    \vspace{3.2cm}}}
  \caption{Diagramme de déploiement UML --- Distribution des services Docker Compose}
  \label{fig:diag-deploiement}
\end{figure}

%------------------------------------------------------------
\section{Résultats}
\label{sec:resultats}
%------------------------------------------------------------

Cette section présente les interfaces fonctionnelles de la plateforme
implémentée, commentées au regard des exigences définies au
Chapitre~2. Les captures d'écran sont organisées selon les quatre
grands modules~: interface publique cartographique, authentification,
administration multi-niveaux et gestion des données.

%-----
\subsection{Interface Publique Cartographique}
\label{ssec:res-carte}

L'interface publique est accessible sans authentification depuis la racine
de la plateforme. Elle répond directement aux exigences EF03, EF04, EF05
et EF16 relatives à la cartographie interactive.

La Figure~\ref{fig:res-carte-publique} illustre la carte au niveau de zoom
national~: les 553~paroisses importées sont regroupées en indicateurs de
cluster numériques selon leur densité géographique. Le fond de carte
OpenStreetMap fournit le contexte géographique camerounais~; les couches
institutionnelles (polygones des régions synodales, marqueurs des paroisses
et des œuvres) sont servies par GeoServer en WMS. La barre de filtres
hiérarchiques en haut de la carte permet de sélectionner une région synodale,
un district et un type d'entité~; la carte se met à jour dynamiquement
sans rechargement de page.

\begin{figure}[H]
  \centering
  % ===== IMAGE 1 : Carte publique avec clustering =====
  \fbox{\parbox{0.88\textwidth}{\centering\vspace{3.8cm}
    \textit{[Capture d'écran 1 --- Carte interactive publique avec clustering de marqueurs]}%
    \vspace{3.8cm}}}
  \caption{Interface publique --- Carte interactive avec regroupement automatique des marqueurs}
  \label{fig:res-carte-publique}
\end{figure}

La Figure~\ref{fig:res-popup} montre le panneau de détail qui s'affiche lors
d'un clic sur un marqueur individuel~: nom de la paroisse, district et région
d'appartenance, coordonnées GPS, contacts et statut d'activité. Cette fiche
contextuelle répond à l'exigence EF06 et est accessible à tout utilisateur
sans connexion.

\begin{figure}[H]
  \centering
  % ===== IMAGE 2 : Popup de détail d'une paroisse =====
  \fbox{\parbox{0.88\textwidth}{\centering\vspace{3.8cm}
    \textit{[Capture d'écran 2 --- Popup de détail d'une paroisse sur la carte]}%
    \vspace{3.8cm}}}
  \caption{Interface publique --- Popup de détail d'une paroisse sélectionnée}
  \label{fig:res-popup}
\end{figure}

%-----
\subsection{Interface d'Authentification}
\label{ssec:res-auth}

La Figure~\ref{fig:res-login} présente la page de connexion de l'espace
d'administration (\texttt{/admin/login}). Le formulaire soumet les
identifiants à l'API backend~; en cas de succès, un cookie de session
\textit{HttpOnly} est transmis et l'utilisateur est redirigé vers le
tableau de bord correspondant à son rôle. En cas d'identifiants incorrects,
un message d'erreur est affiché sans révéler si l'identifiant ou le mot de
passe est la source de l'erreur, conformément aux bonnes pratiques de
sécurité recommandées par l'OWASP~\cite{owasp2021}. Pour un compte dont
l'indicateur de première connexion est actif, la redirection s'effectue
automatiquement vers le formulaire de changement de mot de passe.

\begin{figure}[H]
  \centering
  % ===== IMAGE 3 : Page de connexion administrateur =====
  \fbox{\parbox{0.88\textwidth}{\centering\vspace{3.8cm}
    \textit{[Capture d'écran 3 --- Page de connexion de l'espace administrateur]}%
    \vspace{3.8cm}}}
  \caption{Interface d'authentification --- Page de connexion sécurisée}
  \label{fig:res-login}
\end{figure}

%-----
\subsection{Interface d'Administration Multi-Niveaux}
\label{ssec:res-admin}

L'interface d'administration adapte son contenu et ses fonctionnalités au
rôle de l'utilisateur connecté, conformément aux exigences EF02 et EF07 à
EF09. Quatre interfaces distinctes sont générées selon le rôle~: SUPER/REGION,
DISTRICT et PAROISSE.

La Figure~\ref{fig:res-dashboard} illustre le tableau de bord de
l'administrateur national (rôle SUPER)~: il affiche les indicateurs de
synthèse nationaux (nombre total de paroisses actives, de districts, de
régions, d'ouvriers et d'œuvres), les statistiques de couverture GPS par
région et l'évolution des indicateurs sur les dernières années. Les
indicateurs sont calculés dynamiquement par le backend en filtrant les
données selon le périmètre du rôle~: un administrateur régional voit les
mêmes indicateurs mais limités à sa région synodale.

\begin{figure}[H]
  \centering
  % ===== IMAGE 4 : Tableau de bord administrateur national =====
  \fbox{\parbox{0.88\textwidth}{\centering\vspace{3.8cm}
    \textit{[Capture d'écran 4 --- Tableau de bord de l'administrateur national]}%
    \vspace{3.8cm}}}
  \caption{Interface d'administration --- Tableau de bord de l'administrateur national}
  \label{fig:res-dashboard}
\end{figure}

La Figure~\ref{fig:res-paroisses} présente l'interface de gestion des
paroisses~: liste paginée filtrée par région et district, formulaire de
création et de modification avec validation des coordonnées GPS en temps
réel, et action de suppression avec confirmation. Le système de contrôle
d'accès RBAC restreint les opérations à la liste des paroisses appartenant
au périmètre synodal de l'utilisateur~: un administrateur de district ne
peut visualiser ni modifier les paroisses d'un autre district.

\begin{figure}[H]
  \centering
  % ===== IMAGE 5 : Interface de gestion des paroisses =====
  \fbox{\parbox{0.88\textwidth}{\centering\vspace{3.8cm}
    \textit{[Capture d'écran 5 --- Interface de gestion des paroisses (liste et formulaire)]}%
    \vspace{3.8cm}}}
  \caption{Interface d'administration --- Gestion des paroisses avec contrôle RBAC}
  \label{fig:res-paroisses}
\end{figure}

%-----
\subsection{Module d'Import et d'Export}
\label{ssec:res-import}

Le module d'import (Figure~\ref{fig:res-import}) répond à l'exigence EF10~:
il permet de charger les fichiers Excel fournis par les services de l'EEC
directement dans la base de données, après une validation ligne par ligne.
Le rapport d'import affiché après traitement détaille le nombre de lignes
insérées avec succès, le nombre de lignes rejetées et, pour chaque
rejet, la raison précise (coordonnées GPS hors territoire, district
inexistant, doublon détecté). Cette traçabilité de l'import répond
directement au problème de qualité des données identifié lors de l'audit
préliminaire~: 127~paroisses sources sans coordonnées GPS valides ont ainsi
été identifiées et signalées lors de la première importation.

\begin{figure}[H]
  \centering
  % ===== IMAGE 6 : Module d'import Excel et rapport de résultats =====
  \fbox{\parbox{0.88\textwidth}{\centering\vspace{3.8cm}
    \textit{[Capture d'écran 6 --- Module d'import Excel avec rapport détaillé de résultats]}%
    \vspace{3.8cm}}}
  \caption{Module d'import --- Chargement de données Excel et rapport de validation}
  \label{fig:res-import}
\end{figure}

%-----
\subsection{Journal d'Audit}
\label{ssec:res-audit}

Le journal d'audit (Figure~\ref{fig:res-audit}) est accessible
exclusivement à l'administrateur national. Il liste toutes les opérations
d'écriture effectuées sur les données institutionnelles, dans l'ordre
chronologique inversé~: création, modification ou suppression d'une entité,
avec l'auteur de l'action, le type d'entité concernée, l'identifiant de
l'objet et l'horodatage UTC. Des filtres par type d'action, type d'entité
et période permettent de cibler une recherche. En cliquant sur une entrée,
l'administrateur accède au détail de la modification~: valeurs avant et
après la mise à jour. Ce module répond aux exigences EF14 et EF15 relatives
à la traçabilité, et constitue l'un des apports distinctifs par rapport
à la plateforme GeoEEC~v1.0 qui ne disposait d'aucun mécanisme d'audit.

\begin{figure}[H]
  \centering
  % ===== IMAGE 7 : Journal d'audit =====
  \fbox{\parbox{0.88\textwidth}{\centering\vspace{3.8cm}
    \textit{[Capture d'écran 7 --- Journal d'audit des modifications institutionnelles]}%
    \vspace{3.8cm}}}
  \caption{Journal d'audit --- Historique des opérations d'écriture sur les données}
  \label{fig:res-audit}
\end{figure}

%------------------------------------------------------------
\section{Discussion}
\label{sec:discussion}
%------------------------------------------------------------

\subsection{Adéquation aux Objectifs}
\label{ssec:disc-objectifs}

Les six objectifs formulés en Introduction Générale sont confrontés aux
résultats obtenus. L'objectif~1 (modélisation du schéma de données
hiérarchique) est atteint~: le schéma PostGIS implémenté couvre les
22~régions synodales, 137~districts, 553~paroisses importées, 311~œuvres
et 685~ouvriers, avec toutes les contraintes d'intégrité référentielle et
les champs de géolocalisation. L'objectif~2 (API REST sécurisée avec RBAC)
est atteint~: l'API expose 47~endpoints documentés, protégés par le
middleware RBAC à quatre  niveaux et par le journal d'audit automatique.
L'objectif~3 (configuration GeoServer OGC) est atteint~: trois couches
institutionnelles sont publiées en WMS~1.3.0 et WFS~2.0.0 avec styles SLD.
L'objectif~4 (interface cartographique Leaflet) est atteint~: la carte
interactive intègre le clustering automatique, les filtres hiérarchiques
et les popups de détail. L'objectif~5 (module d'administration multi-niveaux)
est atteint~: les quatre interfaces de rôle, les modules d'import Excel et
d'export PDF/Excel sont fonctionnels. L'objectif~6 (conteneurisation Docker)
est atteint~: la plateforme est entièrement déployable avec une commande
unique sur tout serveur disposant de Docker.

%-----
\subsection{Forces de la Solution}
\label{ssec:disc-forces}

\textbf{Sécurité améliorée par rapport à la version antérieure.}
Le remplacement de l'authentification JWT par des sessions Django avec
cookies \textit{HttpOnly} et \textit{SameSite=Strict} élimine le vecteur
d'attaque XSS qui permettait la capture des jetons transmis dans le stockage
local du navigateur. La protection CSRF systématique sur toutes les
requêtes modifiantes renforce ce dispositif, conformément aux recommandations
de l'OWASP~\cite{owasp2021}.

\textbf{Interopérabilité cartographique.}
La publication des couches géographiques via GeoServer en WMS~1.3.0 et
WFS~2.0.0 rend les données de l'EEC consommables par tout client
cartographique compatible OGC~: QGIS, OpenLayers, ou un futur portail
institutionnel camerounais. Aucune solution concurrente analysée au
Chapitre~1 ne proposait cette combinaison.

\textbf{RBAC aligné sur la hiérarchie réelle de l'EEC.}
Le modèle de contrôle d'accès à quatre niveaux --- national, régional,
district et paroissial --- correspond exactement à la structure synodale
de l'institution~: un administrateur de district peut gérer les paroisses
de son périmètre sans accès aux données des districts voisins. Ce niveau
de granularité était absent de la version antérieure.

\textbf{Traçabilité complète et automatique.}
Tout enregistrement de modification, de création ou de suppression
déclenche automatiquement l'écriture d'une entrée dans le journal d'audit
via les signaux Django, sans intervention du développeur dans chaque vue.
Cette automatisation garantit l'exhaustivité de l'historique, impossible
à assurer manuellement.

\textbf{Déploiement reproductible.}
La configuration Docker Compose en cinq services garantit que l'environnement
d'exécution est identique en développement et en production~: plus aucune
divergence de versions de bibliothèques ou de configuration système ne peut
causer des défaillances à l'installation.

Ces cinq forces confirment la pertinence des décisions architecturales
formalisées au Chapitre~2. Elles positionnent la plateforme de façon
distincte par rapport aux solutions analysées dans l'état de l'art, en
particulier sur les axes de sécurité, d'interopérabilité OGC et de
RBAC institutionnel. La section suivante expose les limites structurelles
qui tempèrent la portée de ces résultats.

%------------------------------------------------------------
\section{Limites et Contraintes}
\label{sec:limites}
%------------------------------------------------------------

Tout travail d'ingénierie comporte des contraintes qui en délimitent la
portée. Une évaluation honnête de ces limites est indispensable pour
apprécier correctement la valeur de la solution et orienter les
développements futurs.

\textbf{Couverture GPS incomplète.}
L'audit des données sources a révélé que 127~paroisses sur 565, soit
22~\% de l'effectif total, ne disposent d'aucune coordonnée GPS valide
dans les fichiers Excel fournis par l'EEC. Ces paroisses sont importées
en base de données mais ne s'affichent pas sur la carte interactive.
La plateforme ne peut donc pas prétendre à une couverture cartographique
complète du réseau ecclésiastique~; son utilité pour les décisions de
planification territoriale nationale reste limitée tant que ce taux
n'est pas réduit par une campagne de collecte terrain.

\textbf{Distance à vol d'oiseau.}
La fonction de calcul d'itinéraire repose sur la formule de Haversine,
qui mesure la distance géodésique en ligne droite entre deux points.
Cette mesure sous-estime systématiquement la durée réelle de trajet
sur un réseau routier camerounais caractérisé par des voies sinueuses
et un relief accidenté~; elle ne tient pas compte des barrières
physiques (fleuves, massifs montagneux, absence de pont).
La durée estimée à 60~km/h de moyenne est donc approximative et ne
peut servir de base à des décisions opérationnelles précises.

\textbf{Absence d'application mobile.}
La plateforme est conçue pour des écrans desktop et tablette
(résolution minimale 1\,024~$\times$~768~px, ENF09). Elle n'est pas
adaptée aux smartphones, qui constituent pourtant le principal outil
numérique des administrateurs paroissiaux en zone rurale camerounaise.
En l'absence d'application mobile dédiée, la collecte terrain des
coordonnées GPS des 127~paroisses non géolocalisées reste sans
outillage numérique adapté.

\textbf{Authentification à facteur unique.}
L'accès à l'espace d'administration repose sur un couple
identifiant/mot de passe sans mécanisme de double authentification
(TOTP ou code temporaire par messagerie). Pour des comptes de niveau
national disposant d'un accès complet à l'ensemble des données de
l'EEC, cette limite expose la plateforme aux risques de compromission
par capture ou réutilisation de mot de passe, en l'absence de second
facteur de vérification.

\textbf{Déploiement mono-instance.}
L'architecture Docker Compose actuelle déploie chaque service en
instance unique sur un seul serveur physique. Elle ne prévoit pas de
mécanisme de réplication ou de bascule automatique en cas de défaillance
matérielle~; la disponibilité de la plateforme est donc directement
dépendante de la fiabilité du serveur institutionnel hôte.

\medskip
\noindent\textbf{Bilan du chapitre.}
Ce chapitre a présenté et justifié les choix technologiques de la
plateforme à travers quatre tableaux comparatifs, exposé les résultats
à travers sept interfaces fonctionnelles commentées, analysé les cinq
forces distinctives de la solution dans la section Discussion, puis
identifié cinq limites structurelles dans cette section. La
\textbf{Conclusion Générale} synthétise l'ensemble de la démarche
et formule les perspectives d'amélioration directement dérivées
de ces limites.

\clearpage

% ============================================================
%  CONCLUSION GÉNÉRALE
% ============================================================
\chapter*{Conclusion Générale}
\addcontentsline{toc}{chapter}{Conclusion Générale}
\markboth{Conclusion Générale}{Conclusion Générale}

\noindent
L'Église Évangélique du Cameroun gérait jusqu'alors ses 565~paroisses,
311~œuvres sociales et 685~ouvriers ecclésiastiques répartis sur
22~régions synodales et 137~districts sans aucun outil numérique
centralisé~: pas de visualisation géographique, pas de contrôle d'accès
formalisé selon les périmètres synodaux, pas de journal d'audit des
modifications, pas de publication cartographique interopérable. Ce travail
a abordé ce problème par la conception et le développement d'une
plateforme web cartographique institutionnelle, en s'appuyant sur
l'hypothèse qu'une architecture à services découplés, entièrement
conteneurisée, combinant un serveur cartographique OGC et un contrôle
d'accès hiérarchique aligné sur la structure synodale, permettrait d'y
répondre de façon complète et déployable dans un contexte institutionnel
à ressources limitées.

Cette hypothèse est validée par les résultats obtenus~: la plateforme
implémentée déploie en une seule commande cinq services
--- Django~5~+~GeoDjango, PostgreSQL~+~PostGIS~3.4, GeoServer~2.26,
Next.js~14 et Redis~7~--- et offre une carte interactive avec clustering
de 553~paroisses, des filtres hiérarchiques par région et district, un
tableau de bord adaptatif par rôle, un module d'import/export des données
institutionnelles, et un journal d'audit automatique et exhaustif de toutes
les opérations d'écriture. L'ensemble des six objectifs formulés en
Introduction Générale a été atteint. Par rapport à la plateforme
GeoEEC~v1.0, la solution apporte une valeur ajoutée précise et documentée
sur cinq axes~: sécurité de l'authentification par sessions \textit{HttpOnly},
RBAC aligné sur les quatre niveaux de la hiérarchie synodale réelle,
publication cartographique conforme aux standards OGC WMS et WFS, traçabilité
automatique par journal d'audit, et reproductibilité de l'environnement
d'exécution par conteneurisation complète.

\medskip
\noindent\textbf{Limites.}
Cinq limites structurelles, analysées en détail dans la
section~\ref{sec:limites} du Chapitre~3, tempèrent la portée de ces
résultats. La couverture cartographique reste incomplète~: 127~paroisses
(22~\%) ne disposent d'aucune coordonnée GPS dans les fichiers sources de
l'EEC et ne s'affichent pas sur la carte, ce qui restreint directement
l'utilité de la plateforme pour les décisions de planification territoriale.
Le calcul d'itinéraire repose sur la distance à vol d'oiseau et non sur
le réseau routier réel, produisant des durées de trajet sous-estimées dans
les zones à relief accidenté. L'absence d'interface mobile prive les
administrateurs paroissiaux d'un outil de saisie GPS sur le terrain.
L'authentification à facteur unique expose les comptes nationaux à des
risques de compromission par mot de passe seul. Enfin, le déploiement
mono-instance ne prévoit pas de mécanisme de haute disponibilité en cas
de défaillance du serveur hôte.

\medskip
\noindent\textbf{Perspectives.}
Ces limites définissent directement les axes de développement prioritaires.
La première perspective est le développement d'une application mobile
légère (React Native ou Progressive Web App) permettant aux administrateurs
paroissiaux de saisir ou de corriger les coordonnées GPS directement sur le
terrain~: cette seule fonctionnalité permettrait de résorber le déficit de
géolocalisation des 127~paroisses non couvertes et de faire passer le taux
de couverture cartographique de 78~\% à 100~\%. La deuxième perspective est
l'intégration d'un moteur de calcul d'itinéraires routiers~--- via l'API
OSRM (\textit{Open Source Routing Machine}) ou une instance locale de
GraphHopper~--- pour remplacer la distance à vol d'oiseau par des temps de
trajet réalistes sur les routes camerounaises. La troisième perspective est
l'activation de la double authentification par TOTP (\textit{Time-based
One-Time Password}) pour les comptes d'administrateurs nationaux et
régionaux, en s'appuyant sur la bibliothèque \texttt{django-otp}
déjà compatible avec l'infrastructure Django. Enfin, la réintégration
du module de questions-réponses en langage naturel développé dans
GeoEEC~v1.0 --- enrichi des nouvelles entités et des statistiques annuelles
--- ouvrirait la plateforme à des usages analytiques avancés permettant
aux responsables synodaux d'interroger les données institutionnelles sans
compétences techniques.

\noindent\rule{\linewidth}{0.8pt}
